% ============================================================================
% LACC: Language-Aware Context Compression
% Nén Ngữ Cảnh Có Nhận Thức Ngôn Ngữ cho Mô Hình Ngôn Ngữ Lớn Tiếng Việt
% ============================================================================
\documentclass[12pt,a4paper]{article}

% ── Encoding & Fonts ──────────────────────────────────────────────────
\usepackage[utf8]{vntex}
\usepackage{fontenc}
\usepackage{newtxtext,newtxmath}
\usepackage{indentfirst}

% ── Layout ────────────────────────────────────────────────────────────
\usepackage[margin=2.5cm]{geometry}
\usepackage{setspace}
\onehalfspacing

% ── Packages ──────────────────────────────────────────────────────────
\usepackage{graphicx}
\usepackage{amsmath,amssymb,amsthm}
\usepackage{algorithm}
\usepackage{algpseudocode}
\usepackage{booktabs}
\usepackage{multirow}
\usepackage{hyperref}
\usepackage{cite}
\usepackage{color}
\usepackage{xcolor}
\usepackage{enumitem}
\usepackage{array}
\usepackage{float}

% ── Theorem environments ──────────────────────────────────────────────
\newtheorem{definition}{Định nghĩa}[section]
\newtheorem{theorem}{Định lý}[section]
\newtheorem{lemma}{Bổ đề}[section]

% ── Custom commands ───────────────────────────────────────────────────
\newcommand{\tone}[1]{\textit{#1}}          % tone name
\newcommand{\score}[1]{\mathcal{S}(#1)}      % scoring function
\newcommand{\weight}[1]{\mathbf{w}_{#1}}     % weight vector
\newcommand{\token}[1]{\texttt{#1}}          % token
\newcommand{\method}[1]{\textsc{#1}}         % method name
\newcommand{\viet}[1]{\textit{#1}}           % Vietnamese word

% ── Title ─────────────────────────────────────────────────────────────
\title{\textbf{LACC: Language-Aware Context Compression \\
        Nén Ngữ Cảnh Có Nhận Thức Ngôn Ngữ \\
        cho Mô Hình Ngôn Ngữ Lớn Tiếng Việt}}

\author{Anonymous ACL submission}
\date{}

\begin{document}
\maketitle

% ============================================================================
\begin{abstract}
Các phương pháp nén ngữ cảnh (\textit{context compression}) hiện tại như LLMLingua, SnapKV hay Selective Context đều được thiết kế theo hướng \textbf{ngôn-ngữ-bất-khả-tri} (\textit{language-agnostic}), và tiếng Việt --- ngôn ngữ đơn lập, có thanh điệu, với token inflation 1.5--2.0$\times$ so với tiếng Anh --- chưa có một benchmark nén ngữ cảnh nào kiểm soát rò dữ liệu (\textit{data leakage}) để đo các phương pháp này một cách đáng tin cậy. Đóng góp chính của bài báo là \method{VCC-Bench} v2 --- benchmark nén ngữ cảnh tiếng Việt với 5 tác vụ (QA văn bản dài, hội thoại đa lượt, needle-in-haystack có kiểm soát nhiễu thanh điệu, gọi công cụ agent, nén xuyên ngôn ngữ), tách bạch tường minh tập huấn luyện/đánh giá để tránh rò dữ liệu.

Trên nền benchmark này, chúng tôi đánh giá có kiểm soát một câu hỏi cụ thể: liệu tín hiệu ngôn ngữ học đặc thù tiếng Việt --- cụ thể là \textbf{thanh điệu} --- có tương quan với vị trí chứa câu trả lời trong tác vụ hỏi-đáp hay không. Chúng tôi đề xuất \method{LACC} (\textit{Language-Aware Context Compression}), một khung nén cấu hình được kết hợp tín hiệu thanh điệu, hình thái từ và perplexity, cùng các biến thể mở rộng: chấm điểm điều kiện hóa theo câu hỏi (\textit{query-conditioning}, so với LongLLMLingua), chọn lọc mức câu, và một kiến trúc thay thế dựa trên encoder phân loại token (kiểu LLMLingua-2). Phát hiện chính của chúng tôi là một \textbf{kết quả âm có kiểm soát}: tín hiệu thanh điệu \textbf{không} tương quan với việc một token có nằm trong câu chứa đáp án hay không cho tác vụ QA tiếng Việt, dù thanh điệu vẫn là một đặc trưng ngôn ngữ học hợp lệ về mặt hình thái-âm vị. Số liệu định lượng đầy đủ trên \method{VCC-Bench} v2, gồm so sánh với LongLLMLingua, LLMLingua-2 và baseline dịch-rồi-nén, được báo cáo ở \S\ref{sec:results}. \method{LACC} và benchmark đều hoạt động được trên phần cứng hạn chế (GPU tiêu dùng 6~GB), và \method{VCC-Bench} v2 cùng mã nguồn được công bố để cộng đồng có thể kiểm chứng và mở rộng kết quả này cho các ngôn ngữ low-resource khác.
\end{abstract}


% ============================================================================
\section{Giới thiệu}
\label{sec:intro}
% ============================================================================

Mô hình ngôn ngữ lớn (LLM) ngày càng được sử dụng để xử lý các ngữ cảnh dài --- từ phân tích tài liệu pháp luật, hội thoại nhiều lượt, đến tác vụ agent đa bước. Tuy nhiên, độ phức tạp bậc hai $O(n^2)$ của self-attention khiến chi phí suy diễn tăng nhanh theo độ dài đầu vào. Để giải quyết vấn đề này, nhiều phương pháp nén ngữ cảnh đã được đề xuất, chia thành ba hướng chính:

\begin{enumerate}[label=(\roman*)]
    \item \textbf{Nén prompt} (LLMLingua~\cite{jiang2023llmlingua}, Selective Context): loại bỏ token/câu không quan trọng dựa trên perplexity hoặc độ tương đồng;
    \item \textbf{Nén KV cache} (SnapKV~\cite{li2024snapkv}, H2O~\cite{zhang2023h2o}, StreamingLLM~\cite{xiao2024streamingllm}): giữ lại KV cache của các token quan trọng dựa trên attention score;
    \item \textbf{Nén không gian tiềm ẩn} (Gist Tokens~\cite{mu2023gist}, ICAE~\cite{ge2024icae}): học các memory token đặc biệt để chứa thông tin nén.
\end{enumerate}

Tất cả các phương pháp trên đều \textbf{mù ngôn ngữ} --- chúng không phân biệt đặc điểm ngôn ngữ của văn bản đầu vào. Với tiếng Anh, điều này ít gây vấn đề vì tiếng Anh không có thanh điệu và có tỉ lệ hư từ thấp hơn. Tuy nhiên, với tiếng Việt --- một ngôn ngữ \textbf{đơn lập, có thanh điệu} (6 thanh) với \textbf{30--40\% token là hư từ} --- việc nén mù ngôn ngữ gây ra ba vấn đề nghiêm trọng:

\begin{enumerate}[label=\textbf{P\arabic*}:,leftmargin=*]
    \item \textbf{Mất thông tin thanh điệu.} Tiếng Việt có 6 thanh điệu tương phản (\tone{ngang}, \tone{huyền}, \tone{sắc}, \tone{hỏi}, \tone{ngã}, \tone{nặng}). Việc xóa token mang dấu thanh có thể thay đổi hoàn toàn nghĩa: \viet{má} (mẹ) $\to$ \viet{ma} (ma), \viet{mả} (mồ mả) $\to$ \viet{ma} (ma). Các phương pháp hiện tại không có cơ chế bảo vệ thông tin thanh điệu.
    
    \item \textbf{Lãng phí ngân sách cho hư từ.} Khoảng 30--40\% token trong văn bản tiếng Việt là hư từ (\viet{đã}, \viet{sẽ}, \viet{của}, \viet{và}, \viet{những}\ldots) --- những từ mang ít thông tin ngữ nghĩa nhưng được giữ lại như thực từ trong các phương pháp nén hiện tại.
    
    \item \textbf{Token inflation.} Cùng một nội dung, tiếng Việt sử dụng nhiều token hơn tiếng Anh 1.5--2.0 lần do đặc điểm hình thái và tokenizer không tối ưu~\cite{lee2026equity, deng2026language}. Điều này khiến context window bị ``phình to'' nhanh hơn, đòi hỏi tỉ lệ nén cao hơn.
\end{enumerate}

\textbf{Đóng góp của bài báo.} Trước P1--P3, câu hỏi tự nhiên là: một tín hiệu ngôn ngữ học hợp lý về mặt lý thuyết (ví dụ thanh điệu) có thực sự \textit{giúp} một compressor giữ được thông tin cần cho tác vụ hạ nguồn (QA) hay không? Trả lời câu hỏi này đòi hỏi một benchmark kiểm soát được rò dữ liệu, chứ không chỉ một phương pháp mới -- đó là lý do đóng góp chính của bài báo là bộ benchmark, không phải bản thân \method{LACC}:

\begin{enumerate}[label=(\arabic*),leftmargin=*]
    \item \textbf{\method{VCC-Bench} v2} (đóng góp chính): benchmark đánh giá nén ngữ cảnh tiếng Việt đầu tiên có kiểm soát rò dữ liệu tường minh, 5 tác vụ (QA văn bản dài, hội thoại đa lượt, needle-in-haystack với thiết kế A/B/C tách bạch nhiễu thanh điệu khỏi tín hiệu thật, gọi công cụ agent, nén xuyên ngôn ngữ), cùng bộ độ đo evidence/attribution (\S\ref{sec:results}) đo trực tiếp việc nén có giữ được câu chứa đáp án hay không, thay vì chỉ đo ROUGE/BLEU gián tiếp.
    \item \textbf{Một kết quả âm có kiểm soát}: trên benchmark này, chúng tôi đo trực tiếp liệu tín hiệu thanh điệu (\S\ref{sec:query_conditioning}--\S\ref{sec:encoder}) có tương quan với vị trí chứa đáp án cho tác vụ QA hay không, tách bạch khỏi các trục khác (điều kiện hóa câu hỏi, cơ chế chấm điểm) để kết luận không bị nhầm lẫn nguyên nhân -- xem Thảo luận (\S\ref{sec:limitations}) cho phát biểu chính xác của phát hiện này và giới hạn của nó.
    \item \textbf{Khung \method{LACC} cấu hình được}: kết hợp tín hiệu thanh điệu, hình thái từ (phân loại token thành hư từ/thực từ/từ láy/từ ghép, mỗi lớp một chiến lược nén) và perplexity, với các biến thể điều kiện hóa theo câu hỏi, chọn lọc mức câu, và một kiến trúc thay thế dựa trên encoder phân loại token -- dùng làm công cụ đo (không phải kết luận đã định trước) cho mục (2), đồng thời vận hành được ở ba mức phần cứng, từ không cần GPU đến GPU tiêu dùng 6~GB.
\end{enumerate}

Mã nguồn, \method{VCC-Bench} v2, và toàn bộ khung thực nghiệm được công bố (liên kết ẩn danh trong quá trình phản biện, xem phần công bố dữ liệu/mã nguồn).

\textbf{Liên hệ với special theme COLING 2027 ("NLP for Linguistics").} Bài báo này dùng phương pháp NLP (benchmark có kiểm soát, đo lường có đối chứng) để trả lời một câu hỏi ngôn ngữ học cụ thể -- thanh điệu tiếng Việt có mang thông tin chức năng cho một tác vụ hạ nguồn hay không -- cho một ngôn ngữ low-resource, đúng tinh thần của special theme.

% ============================================================================
\section{Nghiên cứu liên quan}
% ============================================================================

\subsection{Nén ngữ cảnh cho LLM}

\textbf{Nén prompt.} LLMLingua~\cite{jiang2023llmlingua} sử dụng perplexity-based scoring với budget controller để nén prompt lên đến 20$\times$ mà ít mất hiệu năng. LongLLMLingua~\cite{jiang2023longllmlingua} mở rộng sang chấm điểm điều kiện hóa theo câu hỏi cho ngữ cảnh dài (\S\ref{sec:query_conditioning}). LLMLingua-2~\cite{pan2024llmlingua2} đóng khung nén thành phân loại token nhị phân trên encoder hai chiều thay vì perplexity một chiều (\S\ref{sec:encoder}). Selective Context chọn token dựa trên cosine similarity với query. Các phương pháp này hoạt động trong không gian token rời rạc và không có cơ chế nhận thức ngôn ngữ.

\textbf{Khoảng cách xuyên ngôn ngữ của compressor đã học (\textit{Lost in Compression}~\cite{lukauskas2026lost}).} Nghiên cứu này kiểm toán bốn compressor trích xuất đã học (gồm LLMLingua-2, Kompress-v2) trên mười ngôn ngữ, hơn 250.000 lượt gọi đánh giá, và cho thấy các compressor được giám sát chỉ bằng dữ liệu tiếng Anh (English-supervised) suy giảm hiệu năng đáng kể khi áp dụng ngoài tiếng Anh -- kết luận chính là \textbf{dữ liệu giám sát tiếng Anh không chuyển ngữ (transfer) tốt sang ngôn ngữ khác}, không riêng gì tiếng Việt (tiếng Việt không nằm trong 10 ngôn ngữ được thử). Đây là lý do baseline dịch-rồi-nén (\textit{translate-then-compress}, VN$\to$EN$\to$LLMLingua/-2, B2 ở Bảng~\ref{tab:baselines}\footnote{Baseline này thắng nén bản địa ở khoảng một nửa chi phí token trong 3/5 ngôn ngữ mà~\cite{lukauskas2026lost} thử -- nếu nó cũng thắng trên tiếng Việt, đó là bằng chứng \textit{ủng hộ} phát hiện của họ, không phải phản bác \method{LACC}; xem \S\ref{sec:limitations}.}) là baseline bắt buộc trong bài này (\S\ref{sec:results}): không đo nó thì tiền đề "tiếng Việt cần một compressor bản địa" chưa được kiểm chứng.

\textbf{Vì sao đây không phải cùng một trục với đóng góp của bài báo.} \cite{lukauskas2026lost} đo \textit{compressor học từ dữ liệu tiếng Anh có transfer được sang ngôn ngữ khác không} -- một câu hỏi về giám sát huấn luyện (\textit{supervision}). Bài báo này đo một câu hỏi khác: \textit{tín hiệu ngôn ngữ học đặc thù tiếng Việt (thanh điệu, hình thái từ) có tương quan với vị trí chứa đáp án cho tác vụ QA hay không} -- một câu hỏi về \textit{cơ chế chấm điểm} (\textit{mechanism}), độc lập với compressor được huấn luyện ở đâu. Hai trục bổ sung cho nhau chứ không trùng lặp: \cite{lukauskas2026lost} nói "đừng tin compressor học tiếng Anh sẽ transfer"; bài báo này nói thêm "và một tín hiệu 'đặc thù ngôn ngữ' nghe hợp lý về mặt ngôn ngữ học cũng không tự động là tín hiệu hữu ích cho QA, trừ khi được đo trên benchmark có kiểm soát rò dữ liệu" -- đóng góp benchmark (\method{VCC-Bench} v2) chính là công cụ để đo điều thứ hai, thứ mà~\cite{lukauskas2026lost} không cung cấp cho tiếng Việt.

\textbf{Nén KV cache.} SnapKV~\cite{li2024snapkv} quan sát attention pattern từ ``observation window'' ở cuối prompt để chọn KV position quan trọng. H2O~\cite{zhang2023h2o} giữ ``heavy hitters'' --- token có cumulative attention score cao. StreamingLLM~\cite{xiao2024streamingllm} phát hiện ``attention sink'' và giữ các token đầu + token gần đây. PyramidKV phân bổ ngân sách KV cache theo hình kim tự tháp giữa các layer. Tất cả đều dựa trên attention score --- một tín hiệu không phản ánh đặc điểm ngôn ngữ.

\textbf{Nén không gian tiềm ẩn.} Gist Tokens~\cite{mu2023gist} huấn luyện model nén prompt vào các ``gist token'' đặc biệt (26$\times$ compression). ICAE~\cite{ge2024icae} sử dụng kiến trúc autoencoder với memory slots. Các phương pháp này cần huấn luyện thêm và phụ thuộc vào kiến trúc model cụ thể.

\subsection{Ngôn ngữ học tiếng Việt và LLM}

Tiếng Việt thuộc ngữ hệ Nam Á, là ngôn ngữ \textbf{đơn lập} (isolating) và \textbf{có thanh điệu} (tonal) với 6 thanh~\cite{nguyen1997vietnamese}. Các đặc điểm này tạo ra thách thức riêng cho LLM:

\begin{itemize}
    \item \textbf{Token inflation:} Nghiên cứu của Lee et al.~\cite{lee2026equity} trên 11 ngôn ngữ Đông Nam Á cho thấy Parity-aware BPE là tokenizer công bằng nhất, nhưng token inflation vẫn ở mức 1.5--2.0$\times$ so với tiếng Anh.
    \item \textbf{Energy divide:} Deng et al.~\cite{deng2026language} chỉ ra rằng energy cho inference khác nhau đến 179$\times$ giữa các ngôn ngữ, với low-resource languages chịu ``double penalty'' --- vừa tốn energy hơn vừa accuracy thấp hơn.
    \item \textbf{Cross-lingual token arbitrage:} Colak~\cite{colak2026cross} đề xuất dịch non-English prompt sang English trước khi nén, đạt 34--47\% token reduction. Tuy nhiên, dịch thuật có thể làm mất sắc thái ngôn ngữ gốc.
    \item \textbf{Tokenizer equity:} Guo et al.~\cite{guo2026brain} chỉ ra tokenization fertility chiếm $\sim$60\% sự khác biệt cross-linguistic trong LLM.
\end{itemize}

\textbf{Khoảng trống nghiên cứu:} Chưa có công trình nào thiết kế phương pháp nén ngữ cảnh có nhận thức về đặc điểm ngôn ngữ cho tiếng Việt nói riêng hay các ngôn ngữ low-resource nói chung.


% ============================================================================
\section{Phương pháp đề xuất}
% ============================================================================

\subsection{Tổng quan}

\method{LACC} (\textit{Language-Aware Context Compression}) là khung nén ngữ cảnh kết hợp ba tín hiệu để tính điểm quan trọng cho từng token:

\begin{equation}\label{eq:combined}
    \score{t} = w_{\text{ppl}} \cdot \tilde{\score{}_{\text{ppl}}}(t) \;+\; w_{\text{tone}} \cdot \tilde{\score{}_{\text{tone}}}(t) \;+\; w_{\text{morph}} \cdot \tilde{\score{}_{\text{morph}}}(t)
\end{equation}

trong đó $\tilde{\score{}}(\cdot)$ là các điểm đã được chuẩn hóa về $[0, 1]$, và $w_{\text{ppl}} + w_{\text{tone}} + w_{\text{morph}} = 1$. Khi không có mô hình ngoài, $w_{\text{ppl}} = 0$. Đây là công thức \textbf{wave-1} (\texttt{morph\_combine='add'} trong code), dùng làm cấu hình mặc định và làm điểm so sánh cho các biến thể wave-2 dưới đây.

\paragraph{Biến thể E2 -- kết hợp nhân (\textit{multiplicative blend}).} Phân tích nội bộ wave-1 cho thấy tổng có trọng số làm loãng tín hiệu perplexity mạnh khi cộng với hai tín hiệu yếu hơn (điểm chất lượng riêng-perplexity vượt điểm kết hợp; xem Thảo luận). \texttt{morph\_combine='multiply'} thay vào đó dùng hình thái từ như một \textit{hệ số điều chỉnh} trên nền perplexity+thanh điệu, không phải một số hạng cộng:
\begin{equation}\label{eq:combined_mult}
    \score{t} = \Big[w_{\text{ppl}}' \cdot \tilde{\score{}}_{\text{ppl}}(t) + w_{\text{tone}}' \cdot \tilde{\score{}}_{\text{tone}}(t)\Big] \times \underbrace{\big(0.5 + \tilde{\score{}}_{\text{morph}}(t)\big)}_{\text{hệ số } \in [0.5,\,1.5]}
\end{equation}
với $w_{\text{ppl}}', w_{\text{tone}}'$ là trọng số gốc của \eqref{eq:combined} (không tái chuẩn hóa). Hình thái từ giờ chỉ có thể nhân điểm nền lên tối đa $1.5\times$ hoặc xuống tối thiểu $0.5\times$, không thể lấn át thứ hạng do perplexity/thanh điệu tạo ra như phép cộng có thể.

Hình~\ref{fig:architecture} minh họa kiến trúc tổng thể của \method{LACC} (dạng cộng, wave-1); các biến thể wave-2 thay đổi bước ``Kết hợp'' và/hoặc bước ``Chọn lọc'' như mô tả trong các mục tiếp theo.

\begin{figure}[H]
\centering
\setlength{\fboxsep}{10pt}
\fbox{%
\begin{minipage}{0.9\textwidth}
\centering
\textbf{Kiến trúc LACC}

\vspace{4pt}
\begin{tabular}{c}
$\underbrace{\text{Context } T \text{ ($n$ tokens)}}_{\text{đầu vào}}$ \\[4pt]
$\downarrow$ \\[2pt]
\begin{tabular}{ccc}
\textbf{Phân tích Thanh điệu} & \textbf{Phân tích Hình thái} & \textbf{Phân tích Perplexity} \\
(CPU, từ điển, $O(n)$) & (CPU, từ điển, $O(n)$) & (GPU/CPU, mô hình 135M, $O(n)$) \\
$\downarrow$ & $\downarrow$ & $\downarrow$ \\
$\score{}_{\text{tone}}(t)$ & $\score{}_{\text{morph}}(t)$ & $\score{}_{\text{ppl}}(t)$
\end{tabular} \\[4pt]
$\downarrow$ \\[2pt]
\textbf{Kết hợp}: $\score{t} = \sum w_i \cdot \tilde{\score{}}_i(t)$ \\[4pt]
$\downarrow$ \\[2pt]
\textbf{Chọn lọc}: $\text{TopK}(\score{\cdot}, n/R)$ \\[4pt]
$\downarrow$ \\[2pt]
$\underbrace{\text{Context } T' \text{ ($n/R$ tokens)}}_{\text{đầu ra}}$
\end{tabular}
\end{minipage}
}
\caption{Kiến trúc tổng thể của \method{LACC}. Ba tín hiệu được tính toán song song, sau đó kết hợp với trọng số và chọn lọc token.}
\label{fig:architecture}
\end{figure}

\subsection{Tín hiệu Thanh điệu (\textit{Tone-Aware Scoring})}

\subsubsection{Mật độ và đa dạng thanh điệu}

Với mỗi token $t$ (được giải mã thành chuỗi ký tự), ta định nghĩa:

\begin{definition}[Mật độ thanh điệu]
$$\rho(t) = \frac{|\{c \in t : \text{tone}(c) \neq \text{\tone{ngang}}\}|}{|t|}$$
trong đó $|t|$ là độ dài token (số ký tự) và $\text{tone}(c)$ là thanh điệu của ký tự $c$.
\end{definition}

\begin{definition}[Đa dạng thanh điệu]
$$\nu(t) = \left|\{\text{tone}(c) : c \in t \land \text{tone}(c) \neq \text{\tone{ngang}}\}\right|$$
\end{definition}

\subsubsection{Trọng số bảo toàn thanh điệu}

\begin{equation}\label{eq:tone_weight}
    \weight{tone}(t) = 1.0 + \alpha \cdot \rho(t) \cdot \left(1 + \beta \cdot \frac{\nu(t)}{6}\right)
\end{equation}

với $\alpha = 0.5$ (mức độ quan trọng cơ bản của thanh điệu) và $\beta = 0.3$ (phần thưởng cho đa dạng thanh điệu). Giá trị $\weight{tone}(t) \in [1.0, 1.65]$. Token càng nhiều ký tự có dấu và càng đa dạng thanh điệu thì trọng số càng cao --- \textbf{càng nên được giữ lại}.

\subsubsection{Hệ số tương phản thanh điệu}

Thanh điệu của một token càng khác biệt với các token lân cận thì token đó càng quan trọng (đánh dấu ranh giới ngữ nghĩa). Ta định nghĩa ma trận tương phản $D_{\text{tone}} \in [0,1]^{6 \times 6}$ dựa trên đặc trưng ngữ âm (register, contour):

\begin{table}[H]
\centering
\caption{Ma trận tương phản thanh điệu $D_{\text{tone}}(a, b)$. Giá trị cao = hai thanh khác biệt về mặt ngữ âm.}
\label{tab:tone_matrix}
\begin{tabular}{l|cccccc}
\toprule
 & \tone{ngang} & \tone{huyền} & \tone{sắc} & \tone{hỏi} & \tone{ngã} & \tone{nặng} \\
\midrule
\tone{ngang} & 0.0 & 0.5 & 0.7 & 0.8 & 0.9 & 0.6 \\
\tone{huyền} & 0.5 & 0.0 & 0.9 & 0.6 & 0.8 & 0.4 \\
\tone{sắc}   & 0.7 & 0.9 & 0.0 & 0.7 & 0.4 & 0.8 \\
\tone{hỏi}   & 0.8 & 0.6 & 0.7 & 0.0 & 0.7 & 0.8 \\
\tone{ngã}   & 0.9 & 0.8 & 0.4 & 0.7 & 0.0 & 0.9 \\
\tone{nặng}  & 0.6 & 0.4 & 0.8 & 0.8 & 0.9 & 0.0 \\
\bottomrule
\end{tabular}
\end{table}

Hệ số tương phản cho token $t$ với tập lân cận $N(t)$ (cửa sổ $\pm 2$ token):

\begin{equation}\label{eq:contrast}
    f_{\text{contrast}}(t) = 1 + \gamma \cdot \frac{1}{|N(t)|} \sum_{n \in N(t)} D_{\text{tone}}\big(\text{tone}(t),\; \text{tone}(n)\big)
\end{equation}

với $\gamma = 0.4$ là hệ số khuếch đại.

\subsubsection{Điểm thanh điệu tổng hợp}

\begin{equation}\label{eq:tone_score}
    \score{}_{\text{tone}}(t) = \weight{tone}(t) \times f_{\text{contrast}}(t)
\end{equation}

Ví dụ: token \token{má} (thanh \tone{sắc}) đứng cạnh \token{ma} (thanh \tone{ngang}) và \token{mà} (thanh \tone{huyền}) có:
\begin{itemize}
    \item $\rho(\token{má}) = 1/1 = 1.0$
    \item $\nu(\token{má}) = 1$
    \item $\weight{tone}(\token{má}) = 1.0 + 0.5 \times 1.0 \times (1 + 0.3 \times 1/6) = 1.525$
    \item $f_{\text{contrast}}(\token{má}) = 1 + 0.4 \times (0.7 + 0.9)/2 = 1.32$
    \item $\score{}_{\text{tone}}(\token{má}) = 1.525 \times 1.32 = 2.01$
\end{itemize}

Trong khi đó, \token{và} (hư từ, không có thanh điệu rõ rệt) có $\score{}_{\text{tone}} \approx 1.0$. Token \token{má} được ưu tiên giữ lại gấp đôi \token{và}.

\subsection{Tín hiệu Hình thái từ (\textit{Morphology-Aware Scoring})}

\subsubsection{Phân loại từ}

Mỗi token $t$ được gán một lớp hình thái $c(t)$ thông qua tra cứu từ điển:

\begin{definition}[Phân loại hình thái từ]
$$c(t) \in \{\text{FUNC}, \text{CONTENT}, \text{REDUP}, \text{COMPOUND}, \text{OTHER}\}$$
\end{definition}

\begin{table}[H]
\centering
\caption{Bốn lớp từ chính và chiến lược nén tương ứng.}
\label{tab:word_classes}
\begin{tabular}{lcccc}
\toprule
\textbf{Lớp} & \textbf{Đặc điểm} & \textbf{Tỉ lệ} & \textbf{Ví dụ} & \textbf{Hệ số} \\
\midrule
Hư từ (FUNC) & Ít thông tin ngữ nghĩa & $\sim$35\% & \viet{đã, sẽ, của, và} & 0.40 \\
Thực từ (CONTENT) & Mang nghĩa chính & $\sim$50\% & \viet{học, Việt Nam, máy tính} & 1.20 \\
Từ láy (REDUP) & Dư thừa ngữ âm & $\sim$5\% & \viet{xinh xắn, đẹp đẽ} & 0.60 \\
Từ ghép (COMPOUND) & Đơn vị ngữ nghĩa & $\sim$8\% & \viet{hợp tác xã, giáo viên} & 1.50 \\
\bottomrule
\end{tabular}
\end{table}

Phân loại sử dụng từ điển đóng gồm $\sim$200 hư từ tiếng Việt và $\sim$80 cặp từ láy phổ biến, cho phép tra cứu $O(1)$ không cần mô hình.

\subsubsection{Hệ số bảo toàn theo lớp từ}

\begin{equation}\label{eq:class_mult}
    f_{\text{class}}(c) = \begin{cases}
        0.40 & \text{nếu } c = \text{FUNC} \quad \text{(nén mạnh)} \\
        1.20 & \text{nếu } c = \text{CONTENT} \quad \text{(bảo vệ)} \\
        0.60 & \text{nếu } c = \text{REDUP} \quad \text{(gộp cặp)} \\
        1.50 & \text{nếu } c = \text{COMPOUND} \quad \text{(giữ chặt)} \\
        1.00 & \text{nếu } c = \text{OTHER}
    \end{cases}
\end{equation}

\subsubsection{Phân bổ ngân sách theo lớp}

Thay vì chọn token đơn thuần theo điểm số, \method{LACC} phân bổ ngân sách nén $B$ tỉ lệ với số lượng token trong mỗi lớp, điều chỉnh bởi tỉ lệ giữ lại $r_c$:

\begin{equation}\label{eq:budget}
    B_c = B \times \frac{|T_c| \times r_c}{\sum_{c'} |T_{c'}| \times r_{c'}}
\end{equation}

với $|T_c|$ là số token thuộc lớp $c$, $r_{\text{FUNC}} = 0.30$, $r_{\text{CONTENT}} = 0.85$, $r_{\text{REDUP}} = 0.50$, $r_{\text{COMPOUND}} = 0.95$.

\subsubsection{Xử lý từ láy}

Với các cặp từ láy $(t_L, t_R)$ được phát hiện (ví dụ: $\langle$\viet{xinh}, \viet{xắn}$\rangle$):
\begin{itemize}
    \item Điểm của $t_L$ (từ gốc) được nhân với $1.5$ --- token này đại diện cho cả cặp
    \item Điểm của $t_R$ (từ láy) được gán $0.1$ --- gần như chắc chắn bị loại bỏ
\end{itemize}

\subsection{Tín hiệu Perplexity (\textit{External Model Scoring})}

\subsubsection{Mô hình ngoài}

Khi có GPU, \method{LACC} sử dụng một mô hình ngôn ngữ nhỏ (135M--500M tham số) để ước lượng độ quan trọng của token dựa trên perplexity:

\begin{equation}\label{eq:ppl}
    \score{}_{\text{ppl}}(t_i) = -\log P(t_i \mid t_{<i})
\end{equation}

Token càng ``bất ngờ'' đối với mô hình (perplexity cao) thì càng mang nhiều thông tin và càng nên được giữ lại. Đây là cùng nguyên lý với LLMLingua~\cite{jiang2023llmlingua}, nhưng thay vì dùng mô hình lớn (7B+), \method{LACC} dùng mô hình nhỏ (135M) để tiết kiệm VRAM. \eqref{eq:ppl} được tính trên các cửa sổ trượt $W{=}2048$ token, stride $256$ (E3; bản thảo trước dùng $W{=}512$, cửa sổ hẹp hơn khiến nhiều token đầu mỗi cửa sổ bị chấm điểm với ít ngữ cảnh trái hơn cần thiết) -- xem Thuật toán~\ref{alg:lacc}, bước 3.

\subsubsection{Các mô hình được hỗ trợ}

\begin{table}[H]
\centering
\caption{Các mô hình nhỏ dùng cho perplexity scoring.}
\label{tab:tiny_models}
\begin{tabular}{lcccc}
\toprule
\textbf{Mô hình} & \textbf{Tham số} & \textbf{VRAM (INT4)} & \textbf{Tốc độ} & \textbf{Hỗ trợ tiếng Việt} \\
\midrule
SmolLM2-135M-Instruct & 135M & $\sim$0.3 GB & Nhanh & Cơ bản \\
SmolLM2-360M-Instruct & 360M & $\sim$0.7 GB & Nhanh & Cơ bản \\
Qwen2.5-0.5B-Instruct & 500M & $\sim$0.5 GB & Trung bình & Tốt \\
Qwen2.5-1.5B-Instruct & 1.5B & $\sim$1.0 GB & Chậm hơn & Tốt \\
\bottomrule
\end{tabular}
\end{table}

\subsubsection{Perplexity điều kiện hóa theo câu hỏi (E1, query-conditioning)}
\label{sec:query_conditioning}

Công thức~\eqref{eq:ppl} chấm điểm token chỉ dựa trên ngữ cảnh bên trái, \textbf{không nhìn thấy câu hỏi} $q$ -- một token mang đáp án nhưng dễ đoán từ văn cảnh (perplexity thấp) vẫn bị chấm thấp dù nó chính là điều câu hỏi cần. Đây là nguyên nhân chính (nguyên nhân A, xem Thảo luận) khiến \method{LACC} wave-1 xóa mất câu chứa đáp án trong 53/300 mẫu QA dù giữ nguyên thanh điệu. Theo LongLLMLingua~\cite{jiang2023longllmlingua}, \method{LACC} định nghĩa thêm perplexity \textbf{điều kiện hóa}:
\begin{equation}\label{eq:cond_ppl}
    \score{}_{\text{ppl-cond}}(t_i) = -\log P(t_i \mid q,\, t_{<i})
\end{equation}
với câu hỏi $q$ được ghép vào đầu mỗi cửa sổ trượt (Thuật toán~\ref{alg:lacc}, bước 3) trước ngữ cảnh, sao cho mọi token ngữ cảnh đều được chấm điểm với câu hỏi trong tầm nhìn. Điểm \textbf{tương phản} (\textit{contrastive}) là độ sụt perplexity khi có câu hỏi:
\begin{equation}\label{eq:contrastive_ppl}
    \score{}_{\text{ppl-cx}}(t_i) = \score{}_{\text{ppl}}(t_i) - \score{}_{\text{ppl-cond}}(t_i)
\end{equation}
Token càng \textit{trở nên} dễ đoán hơn khi biết câu hỏi thì càng liên quan đến câu hỏi đó, và càng nên được giữ lại -- ngược với \eqref{eq:ppl}, vốn chỉ ưu tiên token ``bất ngờ'' một cách tổng quát, bất kể câu hỏi là gì. \texttt{contrastive\_ppl=True} (arm \texttt{lacc\_ppl\_contrastive}, \texttt{lacc\_cx\_morph}) dùng \eqref{eq:contrastive_ppl} thay cho \eqref{eq:ppl} khi có $q$; nếu không có câu hỏi, mọi biến thể tự động rơi về \eqref{eq:ppl}.

\paragraph{Hệ số tăng cường từ khoá (query boost).} Độc lập với E1, \textbf{mọi} biến thể \method{LACC} khi có câu hỏi còn nhân thêm một hệ số từ-vựng thô lên $\score{t}$ sau bước kết hợp:
\begin{equation}\label{eq:query_boost}
    \text{qboost}(t; q) = \begin{cases} 1.5 & \text{nếu chuỗi bề mặt của } t \text{ xuất hiện nguyên vẹn trong } q \\ 1.0 & \text{ngược lại} \end{cases}
\end{equation}
Đây là bộ lọc trùng từ khoá (không dùng embedding, không cần mô hình), áp dụng \textbf{sau khi} đã tính $\score{t}$ theo \eqref{eq:combined} hoặc \eqref{eq:combined_mult} -- kể cả trên các arm không bật \texttt{contrastive\_ppl}. Đây là điểm khác biệt quan trọng với mô tả wave-1 gốc: câu hỏi không hoàn toàn vắng mặt khỏi việc chấm điểm ở wave-1, nhưng tham gia rất thô (trùng từ, không phải trùng ngữ nghĩa) và chỉ áp \textit{sau} khi thứ hạng đã được quyết định.

\subsection{Tín hiệu Thanh điệu từ Mô hình (\textit{Model-Based Tone Scoring})}

Ngoài tín hiệu thanh điệu dựa trên từ điển (rule-based), \method{LACC} còn hỗ trợ \textbf{model-based tone scoring}: sử dụng chính hidden states của mô hình đã được huấn luyện với \textit{Phonological Consistency Loss} để tính trọng số bảo toàn. Điều này tạo ra một pipeline khép kín:

\begin{equation}\label{eq:model_tone_pipeline}
    \text{Input} \to \text{Model} \to \text{Hidden States} \to \text{Tone Classifier} \to \text{Tone Scores}
\end{equation}

\subsubsection{Phonological Consistency Loss}

Trong quá trình huấn luyện, một auxiliary loss được thêm vào:

\begin{equation}\label{eq:phon_loss}
    \mathcal{L}_{\text{total}} = \mathcal{L}_{\text{LM}} + \lambda_{\text{tone}} \cdot \mathcal{L}_{\text{tone}}
\end{equation}

\begin{equation}\label{eq:tone_loss}
    \mathcal{L}_{\text{tone}} = \frac{1}{N} \sum_{i=1}^{N} \text{CE}\big(\text{MLP}(h_i),\; y_i^{\text{tone}}\big)
\end{equation}

trong đó $h_i$ là hidden state tại vị trí $i$, $y_i^{\text{tone}} \in \{0,\dots,6\}$ là nhãn thanh điệu của token tại vị trí $i$, và MLP là một classifier 3-layer (Linear $\to$ GELU $\to$ LayerNorm $\to$ Linear). \textbf{Quan trọng}: tone labels được tính trên \textit{vị trí token đã nén} (compressed positions), không phải trên toàn bộ văn bản gốc --- đảm bảo hidden states và tone labels có cùng sequence length, tránh lỗi dimension mismatch.

\subsubsection{Tone Preservation Probe}

Để đánh giá chất lượng bảo toàn thanh điệu, \method{LACC} sử dụng một \textbf{Tone Preservation Probe} --- một linear classifier độc lập huấn luyện nhanh (100--200 steps) trên frozen hidden states để đo tone prediction accuracy:

\begin{equation}\label{eq:tpr}
    \text{TPR} = \frac{|\{i : \hat{y}_i^{\text{tone}} = y_i^{\text{tone}}\}|}{|\{i : y_i^{\text{tone}} \neq \text{ignore}\}|}
\end{equation}

TPR càng cao chứng tỏ hidden states càng mã hóa tốt thông tin thanh điệu sau khi nén.

\subsubsection{Inference: Từ Training đến Compression}

Khi inference với \texttt{use\_model\_tone=True}, \method{LACC} không còn dùng rule-based text analysis nữa mà chạy model forward pass, lấy hidden states, và dùng chính \texttt{PhonologicalConsistencyLoss.score\_importance()} để tính tone-aware score:

\begin{equation}\label{eq:model_score}
    \score{}_{\text{tone-model}}(t_i) = 1.0 + \max_k \text{softmax}(\text{MLP}(h_i))_k
\end{equation}

Token có tone prediction confidence cao $\to$ model nắm rõ thanh điệu $\to$ nên giữ lại. Đây là cầu nối trực tiếp giữa training và inference mà các phương pháp trước đây còn thiếu.

\subsubsection{Biến thể mục tiêu probe (E4): Relevance thay Phonological}

\eqref{eq:tone_loss}--\eqref{eq:model_score} huấn luyện probe để dự đoán \textbf{thanh điệu} của từng vị trí -- một mục tiêu ngôn ngữ học đúng đắn nhưng không trực tiếp đo ``vị trí này có cần cho câu trả lời không''. \texttt{RelevanceConsistencyLoss} là một probe \textbf{cùng kiến trúc, cùng giao diện} (chung lớp phân loại 3 tầng, chung \texttt{score\_importance()}) nhưng đổi nhãn giám sát từ thanh điệu $y_i^{\text{tone}} \in \{0,\ldots,6\}$ sang nhãn nhị phân \textit{liên-quan-câu-hỏi} $y_i^{\text{rel}} \in \{0,1\}$ (token có nằm trong answer span hay không, lấy từ \texttt{qa.jsonl}):
\begin{equation}\label{eq:rel_loss}
    \mathcal{L}_{\text{rel}} = \lambda_{\text{rel}} \cdot \frac{1}{N} \sum_{i=1}^{N} \text{CE}_{\mathbf{w}}\big(\text{MLP}(h_i),\; y_i^{\text{rel}}\big)
\end{equation}
với $\text{CE}_{\mathbf{w}}$ là cross-entropy có trọng số lớp (lớp \textit{relevant} chỉ chiếm vài phần trăm một cửa sổ 512 token -- CE không trọng số bị tối thiểu hoá bằng cách gọi mọi token \textit{irrelevant}, cho accuracy $\sim$0.96 nhưng F1 dương gần 0). Ở suy luận, \eqref{eq:model_score} được dùng nguyên vẹn, chỉ thay $y^{\text{tone}}$ bằng $y^{\text{rel}}$: token có xác suất \textit{relevant} cao được ưu tiên giữ.

\eqref{eq:phon_loss}--\eqref{eq:tone_loss} (Phonological) và \eqref{eq:rel_loss} (Relevance) là \textbf{hai mục tiêu probe loại trừ nhau} trên cùng một khe cắm (\texttt{tone\_source='model'}): bài báo trình bày cả hai và so sánh trực tiếp bằng A/B có kiểm soát trên cùng một SLM đã đóng băng, cùng tín hiệu perplexity/hình thái, cùng generation model (\S\ref{sec:results} -- \texttt{scripts/verify\_tone\_probe\_e2e.py}), thay vì khẳng định trước mục tiêu nào tốt hơn cho tác vụ QA.

\subsection{Tín hiệu Attention (\textit{SnapKV-Style Importance})}

\method{LACC} tích hợp tín hiệu attention pattern từ observation window (SnapKV~\cite{li2024snapkv}) vào \texttt{CombinedCompressor}:

\begin{equation}\label{eq:attention_importance}
    \text{attn\_imp}(t) = \frac{1}{L \cdot H \cdot |W|} \sum_{l=1}^{L} \sum_{h=1}^{H} \sum_{w \in W} \text{Attn}^{l,h}_{w, t}
\end{equation}

trong đó $W$ là observation window (64 token cuối prompt), $L$ là số layer, $H$ là số head. Tín hiệu attention được blend với tone và morphology weights qua tham số $w_{\text{attn}}$:

\begin{equation}\label{eq:combined_v2}
    \score{t} = (1 - w_a)\big[w_t \cdot \tilde{\score{}}_{\text{tone}} + (1-w_t) \cdot \tilde{\score{}}_{\text{morph}}\big] + w_a \cdot \tilde{\text{attn\_imp}}(t)
\end{equation}

\subsection{Cải tiến Hình thái học}

\subsubsection{Phát hiện Từ láy bằng Ngữ âm}

Ngoài dictionary matching, \method{LACC} còn phát hiện cặp từ láy qua \textbf{phonetic similarity}: phân tích âm đầu (initial consonant) và vần (rhyme) của hai âm tiết liền kề:

\begin{equation}\label{eq:phonetic_sim}
    \text{sim}(a, b) = 0.4 \cdot \mathbf{1}[\text{init}(a) = \text{init}(b)] + 0.6 \cdot \mathbf{1}[\text{rhyme}(a) = \text{rhyme}(b)]
\end{equation}

Nếu $\text{sim}(a, b) \geq 0.6$, cặp $(a, b)$ được coi là từ láy và được merge.

\subsubsection{Phân loại Từ Hán-Việt (SINO)}

Từ Hán-Việt chiếm tỉ lệ lớn trong văn bản học thuật và pháp luật tiếng Việt. \method{LACC} thêm lớp \texttt{SINO} với hệ số bảo toàn $f_{\text{sino}} = 1.5$ --- ngang với từ ghép, cao hơn thực từ thông thường ($1.2$). Bộ từ điển Hán-Việt gồm $\sim$60 morpheme phổ biến nhất (\viet{quốc}, \viet{gia}, \viet{chính}, \viet{phủ}, \viet{học}, \viet{sinh}\ldots).

\begin{table}[H]
\centering
\caption{Năm lớp từ mở rộng và chiến lược nén.}
\label{tab:word_classes_v2}
\begin{tabular}{lcccc}
\toprule
\textbf{Lớp} & \textbf{Đặc điểm} & \textbf{Tỉ lệ} & \textbf{Hệ số} \\
\midrule
Hư từ (FUNC) & Ít thông tin ngữ nghĩa & $\sim$35\% & 0.40 \\
Thực từ (CONTENT) & Mang nghĩa chính & $\sim$50\% & 1.20 \\
Từ láy (REDUP) & Dư thừa ngữ âm, detect bằng phonetic sim & $\sim$5\% & 0.60 \\
Từ ghép (COMPOUND) & Đơn vị ngữ nghĩa & $\sim$8\% & 1.50 \\
Từ Hán-Việt (SINO) & Học thuật, pháp luật & $\sim$2\% & 1.50 \\
\bottomrule
\end{tabular}
\end{table}

\subsection{Chọn lọc theo đơn vị câu (E5, sentence-level selection)}
\label{sec:sentence_level}

Chọn lọc mặc định (\eqref{eq:budget} hoặc TopK thuần) hoạt động ở mức \textbf{token}: các token của cùng một câu có thể được giữ/loại độc lập với nhau, kể cả khi câu đó chứa đáp án -- nguyên nhân B trong chẩn đoán wave-1 (nén needle giữ được thanh điệu nhưng mất câu chứa mã kích hoạt, do các token lân cận trong câu bị TopK loại rải rác thay vì cả câu bị loại hoặc giữ như một khối). \texttt{selection\_unit='sentence'} (arm \texttt{lacc\_sentence}) chọn lọc ở mức \textbf{câu}: câu được tách bằng ranh giới câu tiếng Việt (dấu câu kết thúc câu hoặc giới hạn cứng $200$ token/câu -- cùng định nghĩa ranh giới dùng cho baseline \textit{sentence\_retrieval}, để một chiến thắng không chỉ là khác biệt cách tách câu), sau đó câu được xếp hạng theo điểm token trung bình bên trong câu và giữ nguyên cả câu theo thứ tự điểm giảm dần cho tới khi đạt ngân sách:
\begin{equation}\label{eq:sentence_select}
    \bar{\score{}}(s) = \frac{1}{|s|} \sum_{t \in s} \score{t}, \qquad
    \text{kept} = \bigcup \Big\{ s : \bar{\score{}}(s) \text{ nằm trong top theo thứ tự giảm dần, tới khi} \sum_{s \in \text{kept}} |s| \geq \lfloor n/R \rfloor \Big\}
\end{equation}
Câu đầu và câu cuối luôn được giữ (cùng vai trò với $k{=}2$ token biên ở chọn lọc mức token). Nếu ngữ cảnh có ít hơn 3 câu, \texttt{lacc\_sentence} rơi về chọn lọc mức token (\eqref{eq:budget}/TopK) vì tách câu không còn ý nghĩa.

\subsection{Thuật toán tổng hợp}

Thuật toán~\ref{alg:lacc} mô tả quy trình nén đầy đủ của \method{LACC}.

\begin{algorithm}[H]
\caption{\method{LACC}: Language-Aware Context Compression}
\label{alg:lacc}
\begin{algorithmic}[1]
\Require Ngữ cảnh $T$ ($n$ token), tỉ lệ nén $R$, trọng số $w_{\text{ppl}}, w_{\text{tone}}, w_{\text{morph}}$
\Ensure Ngữ cảnh đã nén $T'$ ($\lfloor n/R \rfloor$ token)
\State Tokenize $T \to$ \texttt{token\_ids}$[0..n-1]$
\State Decode \texttt{token\_ids} $\to$ \texttt{tokens}$[0..n-1]$ \Comment{chuỗi ký tự đọc được}
\State
\State // \textbf{Bước 1: Phân tích Thanh điệu} (CPU, $O(n)$)
\For{$i \gets 0$ to $n-1$}
    \State $\rho_i \gets \text{mật\_độ\_thanh\_điệu}(\texttt{tokens}[i])$
    \State $\nu_i \gets \text{đa\_dạng\_thanh\_điệu}(\texttt{tokens}[i])$
    \State $\weight{tone}_i \gets 1.0 + \alpha \cdot \rho_i \cdot (1 + \beta \cdot \nu_i / 6)$
    \State $N \gets \texttt{tokens}[\max(0,i-2)..\min(n,i+2)]$ \Comment{lân cận}
    \State $f_{\text{contrast},i} \gets 1 + \gamma \cdot \text{mean}\{D_{\text{tone}}(\text{tone}(t_i), \text{tone}(n)) : n \in N\}$
    \State $\score{}_{\text{tone},i} \gets \weight{tone}_i \times f_{\text{contrast},i}$
\EndFor
\State
\State // \textbf{Bước 2: Phân tích Hình thái} (CPU, $O(n)$)
\For{$i \gets 0$ to $n-1$}
    \State $c_i \gets \text{phân\_loại\_từ}(\texttt{tokens}[i])$ \Comment{tra từ điển, $O(1)$}
    \State $\score{}_{\text{morph},i} \gets f_{\text{class}}(c_i)$ \Comment{theo công thức~\eqref{eq:class_mult}}
\EndFor
\State Xử lý từ láy: $\forall$ cặp $(i,j)$ là từ láy, $\score{}_{\text{morph},i} \gets \score{}_{\text{morph},i} \times 1.5$, $\score{}_{\text{morph},j} \gets 0.1$
\State
\State // \textbf{Bước 3: Perplexity Scoring} (GPU, tùy chọn, $O(n)$)
\If{có mô hình ngoài}
    \For{mỗi cửa sổ $W{=}2048$ token (stride $256$) trong \texttt{token\_ids}} \Comment{E3: 512 $\to$ 2048}
        \State $\score{}_{\text{ppl}}[i..i{+}W] \gets$ \Call{ComputePerplexity}{tiny\_model, chunk} \Comment{có câu hỏi $\to$ \eqref{eq:contrastive_ppl}}
    \EndFor
\Else
    \State $\score{}_{\text{ppl}} \gets \mathbf{1}$ \Comment{trung tính}
\EndIf
\State
\State // \textbf{Bước 4: Chuẩn hóa \& Kết hợp}
\For{mỗi $\score{} \in \{\score{}_{\text{ppl}}, \score{}_{\text{tone}}, \score{}_{\text{morph}}\}$}
    \State $\tilde{\score{}} \gets (\score{} - \min(\score{})) \;/\; (\max(\score{}) - \min(\score{}))$ \Comment{$\to [0,1]$}
\EndFor
\State $\score{} \gets w_{\text{ppl}} \cdot \tilde{\score{}}_{\text{ppl}} + w_{\text{tone}} \cdot \tilde{\score{}}_{\text{tone}} + w_{\text{morph}} \cdot \tilde{\score{}}_{\text{morph}}$ \Comment{hoặc dạng nhân, \eqref{eq:combined_mult} (E2)}
\If{có câu hỏi $q$}
    \State $\score{} \gets \score{} \odot \text{qboost}(q)$ \Comment{hệ số $1.5\times$ theo trùng khoá, \eqref{eq:query_boost} -- áp cho MỌI biến thể, kể cả khi không bật E1}
\EndIf
\State
\State // \textbf{Bước 5: Chọn lọc}
\State $k \gets 2$ \Comment{số token biên luôn giữ}
\State $B \gets \max(\lfloor n/R \rfloor - 2k,\; 0)$ \Comment{ngân sách cho token giữa}
\If{chọn theo câu (E5)}
    \State $\text{kept} \gets$ \Call{SelectSentences}{$\score{}$, $B$} \Comment{\S\ref{sec:sentence_level}}
\ElsIf{chọn theo hạn ngạch lớp từ (E7)}
    \State $\text{kept} \gets$ \Call{SelectClassProportional}{$\score{}$, $B$} \Comment{\eqref{eq:budget}}
\Else
    \State $\text{kept} \gets \text{TopK}(\score{}[k..n{-}k],\; B)$ \Comment{giữ $B$ token có điểm cao nhất}
\EndIf
\State $T' \gets \texttt{token\_ids}[0..k] \;+\; \text{kept} \;+\; \texttt{token\_ids}[n{-}k..n]$ \Comment{giữ nguyên thứ tự}
\State \Return $T'$
\end{algorithmic}
\end{algorithm}

\subsection{Độ phức tạp}

\begin{table}[H]
\centering
\caption{Độ phức tạp của từng thành phần trong \method{LACC}.}
\label{tab:complexity}
\begin{tabular}{lccc}
\toprule
\textbf{Thành phần} & \textbf{Thời gian} & \textbf{VRAM} & \textbf{Thiết bị} \\
\midrule
Phân tích thanh điệu & $O(n)$ & 0 & CPU \\
Phân tích hình thái & $O(n)$ & 0 & CPU \\
Perplexity scoring & $O(n \cdot d^2)$ & 0.3--1.0 GB & GPU/CPU \\
Kết hợp \& chọn lọc & $O(n \log n)$ & 0 & CPU \\
\midrule
\textbf{Tổng (no-model)} & $O(n \log n)$ & \textbf{0 GB} & CPU \\
\textbf{Tổng (external)} & $O(n \cdot d^2)$ & \textbf{0.3 GB} & GPU \\
\bottomrule
\end{tabular}
\end{table}

\subsection{Kiến trúc thay thế: Encoder Token-Classification (E6)}
\label{sec:encoder}

Mọi biến thể ở trên đều chấm điểm bằng một mô hình \textbf{sinh} (causal LM, decode tuần tự, one-directional) rồi TopK/chọn câu bên trên điểm đó. Wave-1 quan sát rằng một scorer sinh 4B tinh chỉnh ($\sim$8 GB trọng số) vẫn thua một scorer 0.5B có sẵn (0.239 so với 0.442) --- \textit{``scorer tốt hơn không đồng nghĩa compressor tốt hơn''}. E6 thay hẳn kiến trúc: đóng khung nén thành \textbf{phân loại token nhị phân giữ/bỏ} trên một encoder hai chiều (bidirectional, kiểu BERT), theo đúng công thức LLMLingua-2~\cite{pan2024llmlingua2}, thay vì chấm điểm bằng perplexity một chiều.

\paragraph{Huấn luyện (distillation ngoại tuyến).} Một mô hình giáo viên (causal LM, mặc định Qwen2.5-0.5B-Instruct) chấm perplexity từng token bằng \eqref{eq:ppl} (cửa sổ trượt $W{=}2048$) trên corpus không cần câu hỏi (E6 \textbf{query-agnostic} -- nén một lần, không phụ thuộc câu hỏi cụ thể, xem \S\ref{sec:e6_positioning}); $1/R$ token điểm cao nhất được gán nhãn \textit{keep}, phần còn lại \textit{drop}. Nhãn (định nghĩa theo ký tự) được chiếu sang tokenization của encoder Việt (PhoBERT/XLM-R) qua ánh xạ span-ký-tự, rồi một \texttt{AutoModelForTokenClassification} 2 lớp được tinh chỉnh để tái tạo nhãn đó (\texttt{scripts/train\_encoder\_compressor.py}). Không cần \texttt{compression.jsonl} (gold trích xuất bằng LLM) -- nhãn tự sinh trực tiếp từ corpus + giáo viên.

\paragraph{Suy luận.} Ngữ cảnh được tokenize lại bằng tokenizer riêng của encoder (khác tokenizer sinh văn bản); xác suất \textit{keep} $P(\text{keep} \mid t)$ theo mỗi token-encoder được chiếu ngược về ký tự rồi gộp (trung bình) lên từng token của tokenizer sinh văn bản qua cùng cầu nối ký tự dùng cho scorer khác-tokenizer (\S\ref{sec:query_conditioning}), sau đó chọn lọc theo ranh giới top-k như các biến thể khác. Với ngữ cảnh dài hơn cửa sổ encoder, các cửa sổ chồng lấp (stride $128$ token-encoder) được dự đoán riêng rồi lấy trung bình phần chồng lấp.

Vì chạy một lượt xuôi hai chiều duy nhất (không giải mã tuần tự qua nhiều cửa sổ như \eqref{eq:ppl}), E6 rẻ hơn đáng kể về độ trễ suy luận so với mọi biến thể dùng scorer sinh -- đây là trục hiệu quả (\textit{efficiency}) mà E6 nhắm tới, tách biệt khỏi trục chất lượng QA mà E1/E4/E5 nhắm tới (\S\ref{sec:e6_positioning}).


% ============================================================================
\section{Thiết kế thực nghiệm}
% ============================================================================

\subsection{VCC-Bench: Vietnamese Context Compression Benchmark}

Chúng tôi xây dựng \method{VCC-Bench} --- bộ benchmark đầu tiên để đánh giá nén ngữ cảnh cho tiếng Việt, bao gồm 5 tác vụ:

\begin{enumerate}[label=\textbf{T\arabic*}:,leftmargin=*]
    \item \textbf{Long-Document QA.} Hỏi-đáp trên văn bản pháp luật và báo chí tiếng Việt dài (2K--32K token). Đánh giá khả năng giữ lại thông tin chi tiết sau khi nén.
    \item \textbf{Multi-turn Conversation.} Tóm tắt hội thoại nhiều lượt (10--50 lượt). Đánh giá khả năng duy trì mạch hội thoại.
    \item \textbf{Needle-in-Haystack (Vietnamese).} Chèn một câu cụ thể vào ngữ cảnh dài, kiểm tra khả năng truy xuất sau khi nén.
    \item \textbf{Agent Tool-Calling.} Tác vụ agent đa bước với function calling tiếng Việt. Đánh giá khả năng bảo toàn cấu trúc hành động.
    \item \textbf{Cross-lingual Compression.} So sánh nén trực tiếp tiếng Việt với chiến lược ``dịch $\to$ nén $\to$ dịch ngược''.
\end{enumerate}

Kết quả ở \S\ref{sec:results} dùng \textbf{\method{VCC-Bench} v2} (414 mẫu), bản kế thừa có kiểm soát rò dữ liệu tường minh của bộ đầu (T1--T5 giữ nguyên; T3 needle-in-haystack được thiết kế lại thành ba nhóm A/B/C tách bạch tín hiệu thanh điệu thật khỏi nhiễu bề mặt, sau khi phát hiện bản v1 có 16/16 token needle nằm ở sàn trọng số thanh điệu -- tức số liệu v1 đo một bộ dò dấu thanh, không phải đánh đổi thanh điệu-với-QA). Tài liệu chứa mã kích hoạt/PII trong T3 là \textbf{giả có kiểm soát}: giá trị được sinh chương trình (không gắn với cá nhân thật) và đánh dấu tường minh bằng trường \texttt{synthetic\_pii=true} trên mọi mẫu needle (120/120).

\subsection{Mô hình và cấu hình}

\begin{table}[H]
\centering
\caption{Cấu hình thực nghiệm.}
\label{tab:experiment_setup}
\begin{tabular}{ll}
\toprule
\textbf{Tham số} & \textbf{Giá trị} \\
\midrule
Mô hình sinh & Qwen2.5-7B-Instruct (INT4) \\
Mô hình scoring (tùy chọn) & SmolLM2-135M-Instruct (INT4) \\
Tokenizer & Qwen2.5 tokenizer \\
Độ dài ngữ cảnh tối đa & 2,048 token \\
Tỉ lệ nén khảo sát & 2$\times$, 4$\times$, 6$\times$, 8$\times$ \\
Tham số thanh điệu & $\alpha = 0.5$, $\beta = 0.3$, $\gamma = 0.4$ \\
Trọng số (có scorer) & $w_{\text{ppl}} = 0.4$, $w_{\text{tone}} = 0.3$, $w_{\text{morph}} = 0.3$ \\
Trọng số (không scorer) & $w_{\text{ppl}} = 0.0$, $w_{\text{tone}} = 0.5$, $w_{\text{morph}} = 0.5$ \\
GPU & NVIDIA T4 16~GB / P100 16~GB \\
\bottomrule
\end{tabular}
\end{table}

\subsection{Phương pháp so sánh}

Bảng~\ref{tab:baselines} liệt kê toàn bộ các arm được sweep trên VCC-Bench v2 (tên arm khớp 1-1 với \texttt{benchmark.py:WAVE2\_ARMS} và \texttt{run\_pipeline.sh:BENCH\_METHODS}, để một bảng số trong \S\ref{sec:results} có thể truy ngược đúng lệnh đã chạy ra nó).

\begin{table}[H]
\centering
\caption{Toàn bộ arm trong sweep chính (VCC-Bench v2, \S\ref{sec:results}). "Trục" đánh dấu arm thuộc trục A (query-awareness), B (cơ chế chấm điểm), hay cả hai; baseline không thuộc trục nào của \method{LACC}.}
\label{tab:baselines}
\begin{tabular}{llp{6.3cm}}
\toprule
\textbf{Arm} & \textbf{Loại} & \textbf{Mô tả} \\
\midrule
\texttt{none} & Upper bound & Giữ nguyên toàn bộ ngữ cảnh \\
\texttt{random} & Lower bound & Chọn token ngẫu nhiên (nhiều seed, SS4 rule 3) \\
\texttt{llmlingua} & Baseline & LLMLingua~\cite{jiang2023llmlingua} gốc, perplexity một chiều, không điều kiện câu hỏi \\
\texttt{llmlingua\_contrastive} & Baseline (trục A) & LongLLMLingua~\cite{jiang2023longllmlingua}: perplexity điều kiện hóa theo câu hỏi \\
\texttt{translate\_then\_compress} & Baseline (B2) & Dịch VN$\to$EN + LLMLingua, xem \cite{lukauskas2026lost} \\
\texttt{translate\_then\_compress\_llmlingua2} & Baseline (B2) & Dịch VN$\to$EN + LLMLingua-2 (encoder), đúng nghĩa "LLMLingua-2" mà~\cite{lukauskas2026lost} thử \\
\texttt{encoder} & Baseline/\method{LACC}-E6 (trục B) & Encoder token-classification kiểu LLMLingua-2~\cite{pan2024llmlingua2}, \S\ref{sec:encoder} \\
\texttt{lacc\_ppl\_contrastive} & \method{LACC}-E1 (trục A) & Chỉ perplexity điều kiện câu hỏi, tone/morph tắt \\
\texttt{lacc\_ppl\_morph} & \method{LACC}-E2 & Perplexity $\times$ hình thái (nhân), tone tắt \\
\texttt{lacc\_cx\_morph} & \method{LACC}-E1+E2 (trục A) & Perplexity điều kiện câu hỏi $\times$ hình thái \\
\texttt{lacc\_sentence} & \method{LACC}-E5 (trục A) & Chọn lọc mức câu, \S\ref{sec:sentence_level} \\
\texttt{lacc\_classprop} & \method{LACC}-E7 & Ngân sách theo lớp từ, eq.~\eqref{eq:budget} \\
\texttt{lacc\_tone\_gated} & \method{LACC}-E8 & Thanh điệu chỉ bật cho tác vụ bề mặt \\
\texttt{lacc\_tone} & \method{LACC} (wave-1) & Cấu hình gốc (ppl+tone+morph, cộng), điểm đối chứng cho các biến thể trên \\
\bottomrule
\end{tabular}
\end{table}

\subsection{Độ đo}

\begin{itemize}
    \item \textbf{Compression Ratio (CR):} $n_{\text{original}} \;/\; n_{\text{compressed}}$ -- báo cả CR \textit{thực đạt} (achieved) lẫn CR \textit{yêu cầu} (requested), trên cả hai tokenizer (PhoBERT và LLM sinh văn bản), vì hai giá trị này lệch nhau và một bảng chỉ báo CR yêu cầu che giấu việc một arm có đạt ngân sách hay không.
    \item \textbf{Token Savings (TS):} $(n_{\text{original}} - n_{\text{compressed}}) \;/\; n_{\text{original}} \times 100\%$
    \item \textbf{ROUGE-L F1 / BLEU / token-F1 / Exact Match (EM):} Chất lượng sinh văn bản so với đáp án tham chiếu (token-F1 dùng tokenization giữ dấu thanh, \S\ref{sec:results}).
    \item \textbf{Needle recall:} Thu hồi ``đáp án kim'' cho tác vụ needle-in-haystack.
    \item \textbf{Tone Preservation Rate (TPR):} Tỉ lệ token có thanh điệu được bảo toàn sau nén.
    \item \textbf{Evidence recall:} Câu chứa đáp án (xấp xỉ bằng token-overlap cao nhất với đáp án tham chiếu) có sống sót sau nén không.
    \item \textbf{Source-span recoverability:} $P(\text{ngữ cảnh đã nén} \Rightarrow \text{câu chứa đáp án gốc})$, chấm bằng mô hình NLI đa ngôn ngữ.
    \item \textbf{Unsupported-claim rate:} Tỉ lệ câu trong câu trả lời sinh ra không được ngữ cảnh đã nén chứng thực (NLI), một chỉ báo hallucination/attribution.
    \item \textbf{Span-recoverability (hard):} Tỉ lệ token trong ngữ cảnh đã nén xuất hiện \textit{nguyên văn} trong ngữ cảnh gốc ($=1-$ \textit{Variation Rate} của LLMLingua-2~\cite{pan2024llmlingua2}). Đây là chỉ báo \textbf{tất định, không cần NLI}: nén trích nguyên văn đạt $\approx 1.0$, còn dịch-rồi-nén/viết-lại tụt xuống dưới. Nó làm \textit{neo} độc lập để phát hiện khi điểm NLI (source-span recoverability) chứng thực một câu mà kiểm tra chuỗi bác bỏ -- một biện pháp chống \textit{circularity} của việc dùng cùng một mô hình NLI vừa để lọc vừa để chấm.
    \item \textbf{E4 (probe liên-quan):} PR-AUC và recall@budget (2$\times$/4$\times$/8$\times$) thay cho F1 ở ngưỡng cố định (\S\ref{sec:results}).
    \item \textbf{Harmonized Score:} $2 \times Q \times E \;/\; (Q + E)$ với $Q$ = chất lượng, $E$ = hiệu quả.
\end{itemize}

\paragraph{Xử lý thống kê cho kết quả âm.} Vì đóng góp trung tâm là một \textit{kết quả âm}, chúng tôi không dừng ở ``khác biệt không có ý nghĩa'' ($p>0.05$), vốn không chứng minh được sự tương đương. Với mỗi so sánh chính (arm bật thanh điệu so với arm language-agnostic tương ứng), chúng tôi báo cáo: (i) khoảng tin cậy bootstrap ghép cặp cho hiệu số; (ii) \textbf{kiểm định tương đương TOST}~\cite{lakens2017equivalence} với một SESOI (ngưỡng hiệu ứng nhỏ nhất đáng quan tâm) \textbf{đăng ký trước} là 1 điểm EM/F1, để có thể phát biểu ``tương đương trong $\pm 1$ điểm EM'' thay vì chỉ ``không thấy khác biệt''; và (iii) \textbf{phân tích power}~\cite{card2020power}. Ở những tập con nhỏ (ví dụ tác vụ xuyên ngôn ngữ, $n{=}18$) chưa đủ power để khẳng định tương đương, chúng tôi nói rõ điều đó và hạ cấp phát biểu thành ``chưa đủ bằng chứng, cần $n$ lớn hơn'' thay vì diễn giải quá mức.


% ============================================================================
\section{Kết quả thực nghiệm}
\label{sec:results}
% ============================================================================

\subsection{Thử nghiệm khả thi trên GPU phổ thông}
\label{sec:slm_pilot}

Chúng tôi thực hiện một thử nghiệm pilot nhằm kiểm chứng tính khả thi của thành phần mô hình scoring học được trên phần cứng hạn chế. Khác với cấu hình sinh văn bản 7B trong Bảng~\ref{tab:experiment_setup}, thử nghiệm này dùng \texttt{chronopt-research/vietnamese-gpt2-base}, một causal language model tiếng Việt gồm 125.6M tham số. Mô hình được tinh chỉnh bằng LoRA; chỉ 1,179,648 tham số (0.9391\% tổng số tham số) được cập nhật. Loss huấn luyện là tổng của causal-LM loss và \textit{Phonological Consistency Loss}. Tone probe được tối ưu cùng với LoRA adapter.

Thử nghiệm chạy trên NVIDIA GeForce GTX 1060 6~GB với batch size 1, độ dài chuỗi 128 token, LoRA rank 8, gradient accumulation 8 và mixed precision FP16 autocast. Trong bài kiểm tra nhanh, mô hình hoàn thành 30 bước cập nhật; VRAM cực đại do PyTorch ghi nhận là 0.66~GB. Con số này xác nhận cấu hình đủ nhẹ cho GPU 6~GB, nhưng không phải phép đo VRAM toàn hệ thống và có thể thay đổi theo phiên bản CUDA/PyTorch, độ dài chuỗi và batch size.

Bảng~\ref{tab:slm_pilot} báo cáo đánh giá trên 40 văn bản hold-out, được tách với seed cố định 42 từ tập văn bản dùng trong thử nghiệm. Perplexity validation là 51.57 (NLL 3.9430). Tone accuracy trên toàn bộ token là 39.85\%; khi chỉ tính các token được gán thanh khác \tone{ngang}, accuracy là 16.59\% trên 3,141 token. Đây chỉ là kết quả pilot sau 30 bước, chưa phải kết quả hội tụ và chưa có đối chứng base model trên cùng tập validation. Vì vậy, chúng tôi chưa sử dụng kết quả này để khẳng định cải thiện perplexity, chất lượng nén hay TPR; nó chủ yếu chứng minh pipeline LoRA + tone probe có thể huấn luyện, lưu và đánh giá được trên GPU phổ thông.

\begin{table}[H]
\centering
\caption{Kết quả pilot của mô hình scoring tiếng Việt nhỏ sau 30 bước cập nhật (GPU phổ thông). Mũi tên $\downarrow$/$\uparrow$ biểu thị chiều tối ưu.}
\label{tab:slm_pilot}
\begin{tabular}{ll}
\toprule
\textbf{Hạng mục} & \textbf{Giá trị} \\
\midrule
Backbone & \texttt{chronopt-research/vietnamese-gpt2-base} \\
Tổng / tham số LoRA trainable & 125,619,456 / 1,179,648 (0.9391\%) \\
Phần cứng & NVIDIA GeForce GTX 1060 6~GB \\
Cấu hình & batch 1; sequence 128; LoRA $r=8$; accumulation 8 \\
Số bước cập nhật & 30 \\
VRAM PyTorch cực đại & 0.66~GB \\
Số văn bản validation & 40 \\
Validation NLL $\downarrow$ & 3.9430 \\
Perplexity $\downarrow$ & 51.57 \\
Tone accuracy, toàn bộ token $\uparrow$ & 39.85\% \\
Tone accuracy, token có thanh $\uparrow$ & 16.59\% (3,141 token) \\
\bottomrule
\end{tabular}
\end{table}

\paragraph{Cập nhật wave-2: huấn luyện hội tụ trên toàn corpus.} Bảng~\ref{tab:slm_pilot} ở trên là một phép thử khả thi 30 bước, không phải kết quả hội tụ -- \textbf{không dùng số của nó để so sánh chất lượng}. Sau khi hạ tầng dữ liệu \texttt{vncompress-vi-v2} sẵn sàng (\S\ref{sec:query_conditioning}), cùng adapter được huấn luyện đầy đủ \textbf{3 epoch trên toàn bộ 56,729 văn bản / 3,220 tài liệu} (2,394 bước tối ưu, $\sim$7 phút, 1$\times$ H100 80GB -- không còn trên GPU 6~GB, vì đây là phép đo hội tụ chứ không phải phép đo khả thi phần cứng), đánh giá trên tập hold-out 5,673 văn bản / 792,348 token đã chấm điểm:
\begin{center}
\begin{tabular}{ll}
\toprule
\textbf{Hạng mục} & \textbf{Giá trị} \\
\midrule
Validation NLL $\downarrow$ / Perplexity $\downarrow$ & 4.4585 / 86.36 \\
Tone accuracy, toàn bộ token $\uparrow$ & 96.52\% \\
Tone accuracy, token có thanh $\uparrow$ & \textbf{94.85\%} (thay cho 16.59\% ở Bảng~\ref{tab:slm_pilot}) \\
Macro-F1, token có thanh $\uparrow$ & 0.9482 \\
Đối chứng: baseline lớp đa số ("ngang") & 43.48\% \\
Đối chứng: tra bảng token-id không cần huấn luyện & 100.00\% \\
\bottomrule
\end{tabular}
\end{center}
Đối chứng tra-bảng-token-id 100\% cho thấy trần lý thuyết của phép đo này là tầm thường (thanh điệu phần lớn suy ra được trực tiếp từ token id); 94.85\% nói rằng hidden state sau LoRA vẫn giữ được gần hết tín hiệu đó sau fine-tune, không hơn. \textbf{Đây là bằng chứng pipeline huấn luyện + probe hoạt động đúng và hội tụ, KHÔNG PHẢI bằng chứng \method{LACC} nén tốt hơn} -- kết luận đó chỉ đến từ VCC-Bench v2 (\S\ref{sec:results}), theo đúng tinh thần thận trọng đã nêu ở đoạn trên cho bản pilot 30 bước.

\subsection{Hiệu quả nén}

\begin{table}[H]
\centering
\caption{So sánh tỉ lệ nén và thời gian xử lý giữa các phương pháp.}
\label{tab:compression_efficiency}
\begin{tabular}{lcccc}
\toprule
\textbf{Phương pháp} & \textbf{CR thực tế} & \textbf{Tiết kiệm token} & \textbf{Thời gian (ms)} & \textbf{VRAM} \\
\midrule
No Compression & 1.0$\times$ & 0.0\% & --- & 5.0 GB \\
Random Baseline & 4.0$\times$ & 75.0\% & --- & 0 GB \\
LLMLingua (7B) & 4.0$\times$ & --- & --- & 14.0 GB \\
SnapKV & 4.0$\times$ & --- & --- & 14.0 GB \\
LACC-NoModel & 4.0$\times$ & --- & --- & 0 GB \\
LACC-External & 4.0$\times$ & --- & --- & 0.3 GB \\
\bottomrule
\end{tabular}
\end{table}

\subsection{Chất lượng sinh văn bản}

\begin{table}[H]
\centering
\caption{So sánh chất lượng sinh văn bản tại tỉ lệ nén 4$\times$.}
\label{tab:quality}
\begin{tabular}{lcccc}
\toprule
\textbf{Phương pháp} & \textbf{ROUGE-L F1} & \textbf{BLEU} & \textbf{EM} & \textbf{Harmonized} \\
\midrule
No Compression & --- & --- & --- & --- \\
Random Baseline & --- & --- & --- & --- \\
LLMLingua & --- & --- & --- & --- \\
LACC-NoModel & --- & --- & --- & --- \\
LACC-External & --- & --- & --- & --- \\
\bottomrule
\end{tabular}
\end{table}

\subsection{Bảo toàn thanh điệu}

\begin{table}[H]
\centering
\caption{Tỉ lệ bảo toàn thanh điệu (TPR) và tỉ lệ nén hư từ.}
\label{tab:tone_preservation}
\begin{tabular}{lcc}
\toprule
\textbf{Phương pháp} & \textbf{TPR (\%)} & \textbf{Func. Word CR} \\
\midrule
Random Baseline & --- & 4.0$\times$ \\
LLMLingua & --- & --- \\
LACC-NoModel & --- & --- \\
LACC-External & --- & --- \\
\bottomrule
\end{tabular}
\end{table}

\subsection{Phân tích ablation}

\begin{table}[H]
\centering
\caption{Nghiên cứu ablation: đóng góp của từng tín hiệu trong \method{LACC}.}
\label{tab:ablation}
\begin{tabular}{lccc}
\toprule
\textbf{Cấu hình} & \textbf{ROUGE-L} & \textbf{TPR} & \textbf{Func. CR} \\
\midrule
$w_{\text{ppl}}$ only & --- & --- & --- \\
$w_{\text{tone}}$ only & --- & --- & --- \\
$w_{\text{morph}}$ only & --- & --- & --- \\
$w_{\text{tone}} + w_{\text{morph}}$ (no-model) & --- & --- & --- \\
All three (external) & --- & --- & --- \\
\bottomrule
\end{tabular}
\end{table}

\subsection{Ảnh hưởng của tham số}

\begin{figure}[H]
\centering
\fbox{\begin{minipage}{0.8\textwidth}
\centering
\textit{[Hình sẽ được thêm sau:]}\\
(a) Ảnh hưởng của $\alpha$ (tone importance) đến TPR và ROUGE-L.\\
(b) Ảnh hưởng của $w_{\text{tone}}$ (blend weight) đến Harmonized Score.
\end{minipage}}
\caption{Phân tích độ nhạy tham số.}
\label{fig:sensitivity}
\end{figure}


% ============================================================================
\section{Thảo luận}
% ============================================================================

\subsection{Một đặc trưng ngôn ngữ học hợp lý không tự động là tín hiệu hữu ích: trường hợp thanh điệu}
\label{sec:why_tone_negative}

Lập luận P1 (\S\ref{sec:intro}) rằng mất thanh điệu là một rủi ro thực -- \viet{má} và \viet{ma} là hai từ khác nghĩa, và các phương pháp nén mù ngôn ngữ (LLMLingua, SnapKV) không phân biệt hai token này -- \textbf{vẫn đúng về mặt ngôn ngữ học}. Trọng số bảo toàn thanh điệu \eqref{eq:tone_weight} thực sự ưu tiên giữ \token{má} (score $\approx 2.01$) hơn \token{ma} (score $\approx 1.0$), đúng như thiết kế.

Câu hỏi mà benchmark này trả lời là một câu hỏi khác và hẹp hơn: \textbf{với tác vụ QA cụ thể, việc ưu tiên giữ token mang thanh điệu có làm tăng khả năng câu chứa đáp án sống sót sau nén hay không?} Kết quả có kiểm soát trên \method{VCC-Bench} v2 (\S\ref{sec:results}, số liệu điền tại đó) là \textbf{không}: tín hiệu thanh điệu không tương quan với vị trí chứa đáp án. Ba lý do khớp với thiết kế thực nghiệm giải thích vì sao một tín hiệu hợp lý về ngôn ngữ học lại không chuyển hoá thành lợi ích cho QA:

\begin{enumerate}[label=(\arabic*),leftmargin=*]
    \item \textbf{Chấm điểm không biết câu hỏi.} Một token mang thanh điệu đậm nhưng không liên quan tới câu hỏi vẫn được điểm cao như nhau; ngược lại, token chứa đáp án nhưng ít thanh điệu (ví dụ số liệu, tên riêng phiên âm) bị chấm thấp. Đây là lý do trực tiếp dẫn tới E1 (\S\ref{sec:query_conditioning}): tách trục "biết câu hỏi hay không" ra khỏi trục "tín hiệu gì" để không nhầm lẫn hai câu hỏi khác nhau.
    \item \textbf{Đơn vị chọn lọc là token, không phải câu.} Giữ được các token mang thanh điệu rải rác trong một câu không đồng nghĩa với giữ nguyên câu đó -- câu vẫn có thể bị phá vỡ nếu các token lân cận (ít thanh điệu hơn) bị loại. E5 (\S\ref{sec:sentence_level}) kiểm tra trực tiếp giả thuyết này bằng chọn lọc mức câu.
    \item \textbf{Kết hợp cộng làm loãng tín hiệu mạnh hơn.} Khi tín hiệu perplexity (vốn tương quan tốt hơn với thông tin cần cho QA) bị cộng trung bình với tín hiệu thanh điệu, thứ hạng token bị kéo lệch khỏi thứ hạng theo perplexity thuần. E2 (\eqref{eq:combined_mult}) kiểm tra việc kết hợp nhân có giữ được ưu thế của tín hiệu mạnh hơn không.
\end{enumerate}

\textbf{Loại trừ giải thích ``method quá yếu''.} Một kết quả âm chỉ đáng tin khi loại được khả năng pipeline \textit{không bắt được tín hiệu nào} chứ không riêng thanh điệu. Chúng tôi dùng chính arm language-agnostic điều kiện hóa theo câu hỏi (E1, \texttt{lacc\_ppl\_contrastive}, \S\ref{sec:query_conditioning}) làm \textbf{đối chứng dương} (\textit{positive control}): đây là một tín hiệu được kỳ vọng có ích cho QA, chạy qua \textit{cùng} pipeline chấm điểm và \textit{cùng} bộ dữ liệu. Nếu đối chứng dương cải thiện chất lượng QA có ý nghĩa (\S\ref{sec:results}) trong khi thanh điệu thì không, thất bại của thanh điệu là về \textit{tín hiệu}, không phải về \textit{phương pháp} -- đây là điều một probe tương quan đơn lẻ (E4) không tự chứng minh được. Ngoài ra, vì mô hình đọc (\textit{reader}) và encoder PhoBERT vốn đã nhận đầu vào \textit{có dấu}, phép đo của chúng tôi là \textbf{giá trị biên} (\textit{marginal}) của việc \textit{ưu tiên giữ} token mang thanh điệu khi chọn lọc, chứ không phải giá trị của bản thân dấu thanh trong biểu diễn; chúng tôi phát biểu phạm vi hẹp này tường minh để không nhầm ``thanh điệu bị xóa'' với ``thanh điệu không được ưu tiên khi chọn lọc''.

\textbf{Điều phát hiện này KHÔNG nói:} nó không phủ nhận thanh điệu là một đặc trưng ngữ âm-hình thái quan trọng của tiếng Việt, không nói rằng bảo toàn thanh điệu là vô giá trị cho mọi tác vụ (ví dụ dịch máy, phiên âm, trích dẫn nguyên văn -- các tác vụ ``bề mặt'' mà E8/\texttt{lacc\_tone\_gated} vẫn bật tín hiệu này), và không nói rằng không có tín hiệu ngôn ngữ học nào giúp ích cho nén tiếng Việt (hình thái từ, qua phân bổ ngân sách theo lớp E7, là một trục riêng chưa bị bác bỏ bởi phát hiện này). Nó chỉ nói: \textit{đối với QA, thanh điệu như một tín hiệu chấm điểm độc lập-với-câu-hỏi không tương quan với vị trí đáp án} -- một phát biểu hẹp, có kiểm soát, không nên được khái quát hoá thành ``thanh điệu không quan trọng cho tiếng Việt''.

Token inflation (TIR 1.5--2.0$\times$ so với tiếng Anh~\cite{lee2026equity}, nghĩa là context window 128K token chỉ chứa được $\sim$61K từ tiếng Việt so với $\sim$98K từ tiếng Anh) vẫn là một ràng buộc thực của tiếng Việt độc lập với phát hiện trên -- nó là lý do vì sao nén ngữ cảnh cho tiếng Việt vẫn cần thiết, chỉ là tín hiệu "đúng" để làm việc đó chưa chắc là thanh điệu.

\subsection{Định vị E6: trục cơ chế, không phải trục truy vấn}
\label{sec:e6_positioning}

Bài báo phân biệt hai trục cải tiến \textbf{độc lập} với nhau: trục A (\textbf{query-awareness} -- chấm điểm token có nhìn câu hỏi hay không, do E1/E4 phụ trách, \S\ref{sec:query_conditioning}) và trục B (\textbf{cơ chế chấm điểm} -- mô hình sinh một chiều so với encoder phân loại hai chiều, do E6 phụ trách, \S\ref{sec:encoder}). E6 được chốt là \textbf{query-agnostic} (nén một lần, không nhìn câu hỏi cụ thể) có chủ đích: mục tiêu của E6 là cô lập câu hỏi ``kiến trúc chấm điểm nào rẻ/tốt hơn cho cùng một ngân sách nén'', không lẫn với câu hỏi ``có nên điều kiện hóa theo câu hỏi hay không'' mà E1 đã trả lời riêng. Ghép hai trục vào một arm (ví dụ nhân thêm hệ số từ khoá của E1 vào E6) sẽ làm bẩn vai trò \textit{control} này.

Hệ quả trực tiếp cho thiết kế so sánh: E6 được đặt cạnh các baseline \textbf{cùng lớp ngân sách/task-agnostic} -- LLMLingua~\cite{jiang2023llmlingua}, LLMLingua-2~\cite{pan2024llmlingua2}, Selective Context -- chứ \textbf{không} so trực tiếp với LongLLMLingua (query-aware, khác lớp ngân sách: LongLLMLingua được phép dùng câu hỏi để chấm điểm, E6 thì không). Nếu E6 thắng các baseline task-agnostic về chất lượng ở cùng tỉ lệ nén, đó là bằng chứng cho trục B (cơ chế); nếu E6 thua LongLLMLingua, điều đó không phản bác trục B -- nó chỉ xác nhận trục A (query-awareness) quan trọng, một kết luận trục A đã tự đứng vững qua so sánh E1 với LLMLingua thường. Giới hạn (\S\ref{sec:limitations}) nói rõ: E6 không đặt mục tiêu thắng baseline query-aware trên chất lượng QA; mục tiêu của nó là trục hiệu quả (độ trễ, kích thước mô hình).

\subsection{Khả năng mở rộng}

Mặc dù được thiết kế cho tiếng Việt, kiến trúc module hóa của \method{LACC} cho phép mở rộng sang các ngôn ngữ khác:

\begin{itemize}
    \item \textbf{Ngôn ngữ thanh điệu:} Ma trận tương phản $D_{\text{tone}}$ có thể thay thế cho tiếng Trung (4 thanh), tiếng Thái (5 thanh), tiếng Yoruba (3 thanh).
    \item \textbf{Ngôn ngữ đơn lập:} Bộ phân loại hình thái từ có thể mở rộng với từ điển hư từ của tiếng Thái, tiếng Trung, tiếng Miến Điện.
    \item \textbf{Ngôn ngữ low-resource khác:} Tín hiệu perplexity từ mô hình nhỏ hoạt động với bất kỳ ngôn ngữ nào được hỗ trợ bởi tokenizer.
\end{itemize}

\subsection{Hạn chế}
\label{sec:limitations}

\begin{enumerate}[label=(\arabic*),leftmargin=*]
    \item \textbf{Phụ thuộc từ điển:} Tín hiệu thanh điệu và hình thái từ phụ thuộc vào từ điển đóng. Với từ mới hoặc từ mượn, phân loại có thể không chính xác.
    \item \textbf{Mô hình scoring nhỏ:} Perplexity từ mô hình 135M có thể không chính xác bằng mô hình 7B. Tuy nhiên, đây là sự đánh đổi cần thiết cho phần cứng hạn chế.
    \item \textbf{Chưa tối ưu cho streaming:} Thuật toán yêu cầu thấy toàn bộ ngữ cảnh trước khi nén, chưa hỗ trợ streaming từng phần.
    \item \textbf{Trần của E6 (encoder classifier, \S\ref{sec:e6_positioning}):} E6 là \textit{query-agnostic} theo thiết kế, nên không đặt mục tiêu thắng các baseline query-aware (LongLLMLingua, E1/E4) về chất lượng QA -- một compressor không thấy câu hỏi về nguyên tắc không thể chọn đúng bằng một compressor thấy câu hỏi, ở cùng ngân sách. Đóng góp mà E6 nhắm tới là \textit{hiệu quả} (độ trễ suy luận một-lượt-xuôi, kích thước mô hình nhỏ hơn scorer sinh), so với các baseline \textbf{cùng lớp task-agnostic}, không phải chất lượng tuyệt đối so với mọi baseline.
\end{enumerate}

\subsection{Rủi ro}

\begin{enumerate}[label=(\arabic*),leftmargin=*]
    \item \textbf{Khái quát hoá quá mức phát hiện âm.} Kết quả "thanh điệu không tương quan với vị trí đáp án" chỉ được đo cho tác vụ QA trên \method{VCC-Bench} v2, không phải mọi tác vụ hay mọi kiến trúc chấm điểm. Đọc thành "thanh điệu không quan trọng cho NLP tiếng Việt nói chung" là một khái quát hoá sai mà \S\ref{sec:why_tone_negative} cố tình cảnh báo trước.
    \item \textbf{Nén ngữ cảnh trong ứng dụng có rủi ro cao.} Như mọi phương pháp nén ngữ cảnh mất mát (\textit{lossy}), \method{LACC} và các baseline so sánh đều có thể làm mất thông tin cần cho câu trả lời đúng. Triển khai trong lĩnh vực rủi ro cao (y tế, pháp luật) mà không có bước xác minh bổ sung có thể khuếch đại lỗi thay vì chỉ tiết kiệm chi phí suy luận.
    \item \textbf{Nhiễm dữ liệu benchmark (\textit{contamination}).} Vì \method{VCC-Bench} v2 được công bố công khai (\S\ref{sec:results}), các mô hình huấn luyện sau này có thể học thuộc nội dung benchmark, làm mất giá trị của các phép đo held-out trong tương lai -- rủi ro chung của mọi benchmark công khai, không riêng bài này.
\end{enumerate}


% ============================================================================
\section{Kết luận}
% ============================================================================

Bài báo này công bố \method{VCC-Bench} v2 --- benchmark nén ngữ cảnh tiếng Việt có kiểm soát rò dữ liệu, 5 tác vụ, cùng bộ độ đo evidence/attribution đo trực tiếp việc nén có giữ được câu chứa đáp án hay không. Trên nền benchmark này, chúng tôi đo có kiểm soát một tín hiệu ngôn ngữ học cụ thể --- thanh điệu --- và phát hiện nó \textbf{không} tương quan với vị trí chứa đáp án cho tác vụ QA, dù vẫn là một đặc trưng ngữ âm-hình thái hợp lệ của tiếng Việt (\S\ref{sec:why_tone_negative}). Khung \method{LACC} cấu hình được (thanh điệu, hình thái từ, perplexity, cùng các biến thể điều kiện hóa theo câu hỏi, chọn lọc mức câu, và kiến trúc encoder thay thế) là công cụ đo cho phát hiện này, không phải kết luận đã định trước -- vận hành được từ không cần GPU tới GPU tiêu dùng 6~GB.

Hướng phát triển trong tương lai bao gồm: (1) mở rộng \method{VCC-Bench} và phép đo evidence/attribution cho các ngôn ngữ thanh điệu, đơn lập khác; (2) kiểm chứng liệu hình thái từ (E7, chưa bị phát hiện trên bác bỏ) có tương quan với vị trí đáp án tốt hơn thanh điệu hay không; (3) tích hợp với các hệ thống serving như vLLM để nén KV cache; và (4) huấn luyện end-to-end kết hợp \textit{Relevance Consistency Loss} (E4) thay cho \textit{Phonological Consistency Loss}, tuỳ kết quả A/B ở \S\ref{sec:results}.


% ============================================================================
\section{Công bố dữ liệu và mã nguồn}
% ============================================================================

\textbf{Giấy phép.} \method{VCC-Bench} v2 kết hợp ba nguồn với giấy phép khác nhau theo từng tác vụ: CC-BY-SA 4.0 (văn bản Wikipedia, dùng cho T2/T5), Public Domain (văn bản pháp luật, dùng cho T1), và MIT (nội dung tổng hợp cho T3/T4 -- needle, hội thoại, lệnh gọi công cụ). Giấy phép tổng hợp được ghi trong metadata đi kèm bộ dữ liệu.

\textbf{Dữ liệu cá nhân (PII).} Không có dữ liệu cá nhân thật trong \method{VCC-Bench} v2. Nội dung dạng định danh (số tài khoản, số điện thoại, biển số xe...) trong tác vụ T3 (needle-in-haystack) đều là giá trị sinh chương trình, không gắn với cá nhân thật, và được đánh dấu tường minh bằng trường \texttt{synthetic\_pii=true} trên toàn bộ 120/120 mẫu needle -- không cần rà soát thủ công thêm cho trường này.

\textbf{Công bố.} Mã nguồn (\method{LACC}, harness \method{VCC-Bench}, script huấn luyện) và bộ dữ liệu \method{VCC-Bench} v2 sẽ được công bố công khai. Trong giai đoạn phản biện double-blind, liên kết được giữ ẩn danh: \texttt{[LIÊN KẾT ẨN DANH -- BỔ SUNG KHI NỘP]}; liên kết định danh đầy đủ (kèm tên tác giả) sẽ được công bố nếu bài được chấp nhận.


% ============================================================================
% References
% ============================================================================
\bibliographystyle{plain}
\begin{thebibliography}{99}

\bibitem{jiang2023llmlingua}
H. Jiang, Q. Wu, C.-Y. Lin, Y. Yang, and L. Qiu.
\newblock ``LLMLingua: Compressing Prompts for Accelerated Inference of Large Language Models.''
\newblock In \emph{Proceedings of EMNLP}, 2023.

\bibitem{jiang2023longllmlingua}
H. Jiang, Q. Wu, X. Luo, D. Li, C.-Y. Lin, Y. Yang, and L. Qiu.
\newblock ``LongLLMLingua: Accelerating and Enhancing LLMs in Long Context Scenarios via Prompt Compression.''
\newblock \emph{arXiv:2310.06839}, 2023.

\bibitem{pan2024llmlingua2}
Z. Pan, Q. Wu, H. Jiang, M. Xia, X. Luo, J. Zhang, Q. Lin, V. Rühle, Y. Yang, C.-Y. Lin, H. V. Zhao, L. Qiu, and D. Zhang.
\newblock ``LLMLingua-2: Data Distillation for Efficient and Faithful Task-Agnostic Prompt Compression.''
\newblock In \emph{Findings of ACL}, 2024.

\bibitem{li2024snapkv}
Y. Li, Y. Huang, B. Yang, B. Venkitesh, A. Locatelli, H. Ye, T. Cai, P. Lewis, and D. Chen.
\newblock ``SnapKV: LLM Knows What You are Looking for Before Generation.''
\newblock \emph{arXiv:2404.14469}, 2024.

\bibitem{zhang2023h2o}
Z. Zhang, Y. Sheng, T. Zhou, T. Chen, L. Zheng, R. Cai, Z. Song, Y. Tian, C. Ré, C. Barrett, et al.
\newblock ``H2O: Heavy-Hitter Oracle for Efficient Generative Inference of Large Language Models.''
\newblock In \emph{NeurIPS}, 2023.

\bibitem{xiao2024streamingllm}
G. Xiao, Y. Tian, B. Chen, S. Han, and M. Lewis.
\newblock ``Efficient Streaming Language Models with Attention Sinks.''
\newblock In \emph{ICLR}, 2024.

\bibitem{mu2023gist}
J. Mu, X. L. Li, and N. Goodman.
\newblock ``Learning to Compress Prompts with Gist Tokens.''
\newblock In \emph{NeurIPS}, 2023.

\bibitem{ge2024icae}
T. Ge, J. Hu, L. Wang, X. Wang, S.-Q. Chen, and F. Wei.
\newblock ``In-context Autoencoder for Context Compression in a Large Language Model.''
\newblock In \emph{ICLR}, 2024.

\bibitem{lee2026equity}
K. Lee, M. R. Qorib, A. I. Soegeng, and H. T. Ng.
\newblock ``Equity with Efficiency: An Empirical Study of Tokenizers for Multilingual Large Language Models.''
\newblock \emph{arXiv:2606.15044}, 2026.

\bibitem{deng2026language}
N. Deng, A. Shen, Y. Feng, J. Nwatu, J.-W. Chung, M. Chowdhury, Y. Chen, and R. Mihalcea.
\newblock ``The Language-Energy Divide: Measuring Energy Costs of Multilingual LLM Inference.''
\newblock \emph{arXiv:2606.21869}, 2026.

\bibitem{colak2026cross}
M. U. Colak.
\newblock ``Cross-Lingual Token Arbitrage: Optimizing Code Agent Context Windows via Local LLM Preprocessing.''
\newblock \emph{arXiv:2606.03618}, 2026.

\bibitem{lukauskas2026lost}
M. Lukauskas.
\newblock ``Lost in Compression: A Controlled Cross-Lingual Audit of Extractive Prompt Compressors.''
\newblock \emph{arXiv:2608.26175}, 2026.

\bibitem{guo2026brain}
D. Guo, J. Wu, and S. M. Yiu.
\newblock ``Brain-LLM Alignment Tracks Training Data, Not Typology.''
\newblock In \emph{CoNLL}, 2026.

\bibitem{tang2026seco}
J. Tang, Z. Huang, X. Zhang, C. J. Zhang, J. Yu, L. Zheng, R. Meng, and J. Yin.
\newblock ``Beyond Position Bias: Shifting Context Compression from Position-Driven to Semantic-Driven.''
\newblock \emph{arXiv:2605.09463}, 2026.

\bibitem{nguyen1997vietnamese}
D.-H. Nguyen.
\newblock \emph{Tiếng Việt --- Vietnamese.}
\newblock In \emph{London Oriental and African Language Library}, John Benjamins, 1997.

\bibitem{ge2023extractive}
T. Ge, J. Hu, X. Wang, S.-Q. Chen, and F. Wei.
\newblock ``In-context Autoencoder for Context Compression in a Large Language Model.''
\newblock \emph{arXiv:2307.06945}, 2023.

\bibitem{lin2026compresskv}
X. Lin, J. Wang, O. Kondrateva, Y. Shi, B. Li, and G. L. Zhang.
\newblock ``CompressKV: Semantic-Retrieval-Guided KV-Cache Compression for Resource-Efficient Long-Context LLM Inference.''
\newblock \emph{arXiv:2606.24467}, 2026.

\bibitem{kai2026infokv}
J. Kai, Z. Xiao, A. Birch, and Z. Lin.
\newblock ``Information-Aware KV Cache Compression for Long Reasoning.''
\newblock \emph{arXiv:2606.26875}, 2026.

% ── Vietnamese NLP papers ───────────────────────────────────────────

\bibitem{nguyen2020phobert}
D. Q. Nguyen and A. T. Nguyen.
\newblock ``PhoBERT: Pre-trained Language Models for Vietnamese.''
\newblock In \emph{Findings of EMNLP}, 2020.

\bibitem{nguyen2023videberta}
H. S. Nguyen et al.
\newblock ``ViDeBERTa: A Powerful Pre-trained Language Model for Vietnamese.''
\newblock In \emph{EACL}, 2023.

\bibitem{phan2022vit5}
L. Phan et al.
\newblock ``ViT5: Pretrained Transformer-based Models for Vietnamese.''
\newblock \emph{arXiv:2205.06457}, 2022.

\bibitem{vo2024vimistral}
J. Vo.
\newblock ``Vi-Mistral-X: Building a Vietnamese Language Model with Advanced Continual Pre-training.''
\newblock \emph{arXiv:2403.15470}, 2024.

\bibitem{ta2026vidia2std}
K. A. Ta, N. V. Dinh, and K. V. Nguyen.
\newblock ``ViDia2Std: A Parallel Corpus and Methods for Vietnamese Dialect-to-Standard Translation.''
\newblock In \emph{AAAI} (Oral), 2026.

\bibitem{nguyen2018rdrsegmenter}
D. Q. Nguyen et al.
\newblock ``RDRsegmenter: A Fast and Accurate Vietnamese Word Segmenter.''
\newblock In \emph{LREC}, 2018.

\bibitem{nguyen2026vihallu}
A. T.-H. Nguyen et al.
\newblock ``DSC2025 ViHallu Challenge: Detecting Hallucination in Vietnamese LLMs.''
\newblock \emph{arXiv:2601.04711}, 2026.

\bibitem{sophia2025compass}
S. Maria.
\newblock ``Compass-v3: Scaling Domain-Specific LLMs for Multilingual E-Commerce in Southeast Asia.''
\newblock \emph{arXiv:2509.09121}, 2025.

\bibitem{li2023selective}
Y. Li et al.
\newblock ``Selective Context: Compressing Prompts for Efficient In-Context Learning.''
\newblock \emph{arXiv:2310.06201}, 2023.

\bibitem{trukhina2026context}
N. Trukhina and V. Vashkelis.
\newblock ``Context Codec: Compress the Context, Keep the Commitments.''
\newblock \emph{arXiv:2605.17304}, 2026.

\bibitem{ni2026anchorkv}
N. Ni and Y. Lao.
\newblock ``AnchorKV: Safety-Aware KV Cache Compression via Soft Penalty.''
\newblock \emph{arXiv:2606.17872}, 2026.

% ── New Vietnamese NLP & Linguistics papers ────────────────────────

\bibitem{dinh2026hutieubert}
A. T. D. Dinh, K. H. N. Vo, T. T. Ta, V. C. Doan, and T. Quan.
\newblock ``When Morphology Hides in Plain Sight: Breaking the Isolation in Vietnamese and Beyond.''
\newblock In \emph{Proceedings of ACL}, pages 10377--10392, 2026.

\bibitem{nguyen2023phogpt}
D. Q. Nguyen, L. T. Nguyen, C.-K. Tran, D. N. Nguyen, D. Phung, and H. Bui.
\newblock ``PhoGPT: Generative Pre-training for Vietnamese.''
\newblock \emph{arXiv:2311.02945}, 2023.

\bibitem{nguyen2025viwordformer}
N. H. Nguyen, D. T. Nguyen, and N. L.-T. Nguyen.
\newblock ``Vietnamese Words Are Not Constructed from Syllables: Rethinking the Role of Word Segmentation in NLP for Vietnamese Texts.''
\newblock In \emph{Proceedings of AAAI}, 39(22):24069--24077, 2025.

\bibitem{bami2026bamibert}
Qualcomm AI Research.
\newblock ``BamiBERT: A New BERT-based Language Model for Vietnamese.''
\newblock \emph{arXiv:2607.02259}, 2026.

\bibitem{nrl2026diacritic}
NRL-AI.
\newblock ``nrl-ai/vn-diacritic-vit5-base: Vietnamese Diacritic Restoration (ViT5 fine-tune).''
\newblock \emph{HuggingFace Model Hub}, 2026.

\bibitem{nguyen2025vietnormalizer}
H. V. Nguyen, L. Do, T. N. Nguyen, U. S. Khwakhali, T. T. Pham, and V. Do.
\newblock ``VietNormalizer: An Open-Source Library for Vietnamese Text Normalization in TTS and NLP.''
\newblock \emph{arXiv:2603.04145}, 2026.

\bibitem{do2025reference}
T. Do, D. P. Tran, A. Vo, and D. Kim.
\newblock ``Reference-Based Post-OCR Processing with LLM for Precise Diacritic Text in Historical Document Recognition.''
\newblock In \emph{Proceedings of AAAI}, 39(27):27951--27959, 2025.

\bibitem{schmidt2024tokenization}
C. W. Schmidt, V. Reddy, H. Zhang, A. Alameddine, O. Uzan, Y. Pinter, et al.
\newblock ``Tokenization Is More Than Compression.''
\newblock \emph{arXiv:2402.18376}, 2024.

\bibitem{quac2026mutant}
N. Quac et al.
\newblock ``MUTANT: A Recipe for Multilingual Tokenizer Design.''
\newblock In \emph{Proceedings of ACL}, 2026.

\bibitem{huynh2026spelling}
H. T. N. Huynh, L. Nguyen, N. H. Nguyen, H. P. T. Nguyen, and T. T. Quan.
\newblock ``A Two-Stage Vietnamese Spelling Correction Pipeline Combining Underthesea and BARTpho.''
\newblock \emph{VNU Journal of Science: Comp. Science \& Com. Eng.}, 2026.

\bibitem{lakens2017equivalence}
D. Lakens.
\newblock ``Equivalence Tests: A Practical Primer for t Tests, Correlations, and Meta-Analyses.''
\newblock \emph{Social Psychological and Personality Science}, 8(4):355--362, 2017.

\bibitem{card2020power}
D. Card, P. Henderson, U. Khandelwal, R. Jia, K. Mahowald, and D. Jurafsky.
\newblock ``With Little Power Comes Great Responsibility.''
\newblock In \emph{Proceedings of EMNLP}, 2020.

\end{thebibliography}

\end{document}
